% Les fléaux sociaux et leurs conséquences — Allemand 1ère A — Cours signé OnBuch+
% Programme officiel MINESEC. Compiler avec : tectonic ALA03-les-fleaux-sociaux-et-leurs-consequences.tex
\newcommand{\DOCMATIERE}{Allemand}
\newcommand{\DOCNIVEAU}{1ère A}
\newcommand{\DOCMODULE}{Chapitre 1 : Vie familiale et sociale}
\newcommand{\DOCLECON}{ALA03}
\newcommand{\DOCTITRE}{Les fléaux sociaux et leurs conséquences}
% OnBuch+ — préambule « cours signés » ENRICHI (compilé avec tectonic / XeLaTeX)
% Système visuel « Pop » d'OnBuch+ : fond crème (jamais blanc), bordures encre
% épaisses, ombres « dures » décalées, titres Archivo Black en capitales.
% Version MATHÉMATIQUES : graphes (pgfplots/tikz), boîtes pédagogiques
% (définition, propriété, méthode, attention, le savais-tu, exercice résolu,
% exercice, TP, bilan), page de garde et signature OnBuch+ ancrée en bas.
%
% Macros attendues dans main.tex AVANT \input{preamble} :
%   \DOCMATIERE  \DOCNIVEAU  \DOCMODULE  \DOCLECON (numéro)  \DOCTITRE
\documentclass[11pt]{article}

\usepackage{fontspec}
\usepackage[ngerman,french]{babel} % français = langue principale ; l'allemand via \all{..} / textebox
\usepackage[a4paper,margin=2cm,top=3.4cm,bottom=2.8cm,headheight=1.4cm,footskip=1.3cm]{geometry}
\usepackage{amsmath,amssymb,amsfonts}
\usepackage{mathtools} % \xmapsto, \coloneqq... souvent utilisés par les modèles
\usepackage{cancel} % simplifications barrées (bilans de forces)
% \abs{z} (module, valeur absolue) et \norm{u} (norme), utilisés par les modèles
\providecommand{\abs}{}\let\abs\relax\DeclarePairedDelimiter{\abs}{\lvert}{\rvert}
\providecommand{\norm}{}\let\norm\relax\DeclarePairedDelimiter{\norm}{\lVert}{\rVert}
\usepackage[table]{xcolor} % [table] : fournit aussi \rowcolors (alternance de lignes)
\usepackage{pagecolor}
\usepackage{tikz}
\usetikzlibrary{arrows.meta,positioning,calc,shapes.geometric,shapes.misc,shapes.arrows,decorations.pathmorphing,decorations.pathreplacing,decorations.markings,patterns,shadows,mindmap,trees,fit,backgrounds,matrix,intersections,angles,quotes}
\usepackage{pgfplots}
\usepackage[version=4]{mhchem} % équations nucléaires \ce{^{235}_{92}U}
\usepackage[european,siunitx]{circuitikz} % schémas électriques
\pgfplotsset{compat=1.18}
\usepgfplotslibrary{fillbetween,groupplots} % aires entre courbes (calcul intégral)
\usepackage{enumitem}
\usepackage{fancyhdr}
% [most] charge toutes les bibliothèques tcolorbox (listings, minted, raster,
% documentation...) et gonfle inutilement la mémoire TeX sur un document long
% (~300 boîtes) : on ne charge que ce qui est réellement utilisé.
\usepackage[skins,breakable]{tcolorbox}
\usepackage{graphicx}
\usepackage{colortbl}
\usepackage{booktabs}
\usepackage{tabularx}
\usepackage{diagbox} % case d'en-tête diagonale des tableaux à double entrée
% Colonnes extensibles centrée / gauche / droite (C, L, R), souvent utilisées par les modèles
\newcolumntype{C}{>{\centering\arraybackslash}X}
\newcolumntype{L}{>{\raggedright\arraybackslash}X}
\newcolumntype{R}{>{\raggedleft\arraybackslash}X}
\usepackage{array}
\usepackage{multicol}
\usepackage{multirow}
\usepackage{siunitx}
% parse-numbers=false : accepte aussi \SI{2,5 \times 10^{-3}}{...}
\sisetup{locale=FR,output-decimal-marker={,},group-minimum-digits=4,parse-numbers=false}
\DeclareSIUnit\ohm{\ensuremath{\symup{\Omega}}} % U+2126 absent de la police
\DeclareSIUnit\division{div} % réglages d'oscilloscope (ms/div, V/div)
\DeclareSIUnit\tr{tr} % tours (vitesse de rotation, ex. \SI{1500}{\tr\per\minute})
\DeclareSIUnit\molar{M} % concentration molaire (ex. \SI{3}{\molar}, \SI{1}{\micro\molar})
\DeclareSIUnit\an{an} % année (le modèle confond parfois avec \year, un compteur TeX) — ex. \SI{5}{\kilo\meter\per\an}
\DeclareSIUnit\FCFA{FCFA} % franc CFA (ex. \SI{500}{\FCFA\per\kilo\gram})
\DeclareSIUnit\degreeBrix{\textdegree Brix} % teneur en sucre d'un sirop/jus (réfractométrie)
\DeclareSIUnit\month{mois} % le modèle utilise parfois \month, un compteur TeX, comme unité de durée
\usepackage{qrcode}
% \degree : commande fréquemment utilisée par les modèles pour un angle ou une
% température en dehors de \SI{}{} (non fournie par les paquets chargés).
\providecommand{\degree}{^{\circ}}
% \qed (fin de démonstration), utilisé par les modèles sans amsthm
\providecommand{\qed}{\hfill$\square$}
% \barre : double barre des valeurs interdites dans un tableau de variations (tkz-tab)
\providecommand{\barre}{\ensuremath{\|}}
% environnement proof (amsthm non chargé) utilisé par les modèles
\ifdefined\proof\else\newenvironment{proof}[1][Démonstration]{\par\noindent\textit{#1.}\ }{\qed\par}\fi
\usepackage{lastpage}
\usepackage{hyperref}
\hypersetup{hidelinks}

% ============================================================
% Palette « Pop » OnBuch+ — mêmes hex que la classe P (plus_ui.dart)
% ============================================================
\definecolor{poporange}{HTML}{F59321}
\definecolor{popdark}{HTML}{E07A0C}
\definecolor{popcream}{HTML}{FFF6E8}
\definecolor{popink}{HTML}{1C1714}
\definecolor{poppurple}{HTML}{7A5AE0}
\definecolor{popgreen}{HTML}{2E9E6A}
\definecolor{poppink}{HTML}{E0457B}
\definecolor{popblue}{HTML}{2F7DE1}
\definecolor{popgold}{HTML}{C98A12}
\definecolor{popmuted}{HTML}{8A7F76}
\definecolor{popline}{HTML}{EADFD2}
\definecolor{popgray}{HTML}{8A7F76}
\colorlet{popgrey}{popgray}
% Alias tolérants (le modèle nomme parfois les couleurs différemment)
\colorlet{popwhite}{white}
\colorlet{popblack}{popink}
\colorlet{popred}{poppink}
\colorlet{popyellow}{popgold}
% Teintes claires (fonds de tableaux/encadrés) — le modèle nomme parfois
% les couleurs avec un suffixe L (ex. popblueL) sans les avoir définies.
\colorlet{poporangeL}{poporange!20}
\colorlet{popdarkL}{popdark!20}
\colorlet{poppurpleL}{poppurple!20}
\colorlet{popgreenL}{popgreen!20}
\colorlet{poppinkL}{poppink!20}
\colorlet{popblueL}{popblue!20}
\colorlet{popgoldL}{popgold!20}
\colorlet{popredL}{popred!20}
\colorlet{popgrayL}{popgray!20}
% Teintes claires pour fonds de tableaux / zones
\colorlet{popblueL}{popblue!10!white}
\colorlet{poporangeL}{poporange!14!white}
\colorlet{popgreenL}{popgreen!12!white}
\colorlet{poppurpleL}{poppurple!12!white}
\colorlet{poppinkL}{poppink!12!white}
\colorlet{popgoldL}{popgold!14!white}

\pagecolor{popcream}

% ============================================================
% Typographie
% ============================================================
\defaultfontfeatures{Ligatures=TeX}
\setmainfont{PlusJakartaSans-Medium.ttf}[Path=../../../fonts/,BoldFont=PlusJakartaSans-Bold.ttf]
% Symboles, API et emojis absents de la police principale : repli sur DejaVu Sans (caractères actifs)
\newfontface\symfont{DejaVuSans.ttf}[Path=../../../fonts/]
\makeatletter
\newcommand\symactive[1]{\catcode#1=13 \begingroup\lccode`\~=#1 \lowercase{\endgroup\def~}{{\symfont\char#1}}}
\newcommand\symblank[1]{\catcode#1=13 \begingroup\lccode`\~=#1 \lowercase{\endgroup\def~}{}}
\newcommand\symhyph[1]{\catcode#1=13 \begingroup\lccode`\~=#1 \lowercase{\endgroup\def~}{\mbox{-}}}
\newcount\symc
\newcommand\symrange[2]{\symc=#1 \@whilenum\symc<#2 \do{\expandafter\symactive\expandafter{\the\symc}\advance\symc by 1 }}
\newcommand\symblankrange[2]{\symc=#1 \@whilenum\symc<#2 \do{\expandafter\symblank\expandafter{\the\symc}\advance\symc by 1 }}
\makeatother
\symrange{"0250}{"02FF}\symactive{"2713}\symactive{"2714}\symactive{"2717}\symactive{"2718}\symactive{"2610}\symactive{"2611}\symactive{"2612}
\symhyph{"2011}\symhyph{"2E1A}\symblankrange{"1F300}{"1FAFF}\symblank{"FE0F}
% Maths en sans-serif (Fira Math) avec chiffres et lettres droites de Plus
% Jakarta, pour que les formules se fondent dans le texte.
\usepackage[math-style=ISO,bold-style=ISO]{unicode-math}
\setmathfont{FiraMath-Regular.otf}[Path=../../../fonts/]
\setmathfont{PlusJakartaSans-Medium.ttf}[Path=../../../fonts/,range={up/num,up/{Latin,latin},\mathup}]
\usepackage{icomma} % virgule décimale sans espace en mode math
% Caractères mathématiques Unicode absents des polices de texte (Plus Jakarta,
% Archivo Black) : rendus par leur équivalent mathématique, en texte comme en
% formule — sinon ils s'affichent en carré vide (ex. « Arithmétique dans ℤ »).
\usepackage{newunicodechar}
\newunicodechar{ℝ}{\ensuremath{\mathbb{R}}}
\newunicodechar{ℤ}{\ensuremath{\mathbb{Z}}}
\newunicodechar{ℂ}{\ensuremath{\mathbb{C}}}
\newunicodechar{ℕ}{\ensuremath{\mathbb{N}}}
\newunicodechar{ℚ}{\ensuremath{\mathbb{Q}}}
\newunicodechar{𝔻}{\ensuremath{\mathbb{D}}}
\newunicodechar{α}{\ensuremath{\alpha}}
\newunicodechar{β}{\ensuremath{\beta}}
\newunicodechar{γ}{\ensuremath{\gamma}}
\newunicodechar{δ}{\ensuremath{\delta}}
\newunicodechar{ε}{\ensuremath{\varepsilon}}
\newunicodechar{θ}{\ensuremath{\theta}}
\newunicodechar{λ}{\ensuremath{\lambda}}
\newunicodechar{μ}{\ensuremath{\mu}}
\newunicodechar{σ}{\ensuremath{\sigma}}
\newunicodechar{τ}{\ensuremath{\tau}}
\newunicodechar{φ}{\ensuremath{\varphi}}
\newunicodechar{ψ}{\ensuremath{\psi}}
\newunicodechar{ω}{\ensuremath{\omega}}
\newunicodechar{Σ}{\ensuremath{\Sigma}}
% majuscules produites par \MakeUppercase dans les titres (α-aminés -> Α)
\newunicodechar{Α}{\ensuremath{\alpha}}
\newunicodechar{Β}{\ensuremath{\beta}}
\newunicodechar{Γ}{\ensuremath{\Gamma}}
\newunicodechar{Δ}{\ensuremath{\Delta}}
\newunicodechar{Ε}{\ensuremath{\varepsilon}}
\newunicodechar{Θ}{\ensuremath{\Theta}}
\newunicodechar{Λ}{\ensuremath{\Lambda}}
\newunicodechar{Μ}{\ensuremath{\mu}}
\newunicodechar{Τ}{\ensuremath{\tau}}
\newunicodechar{Φ}{\ensuremath{\Phi}}
\newunicodechar{Ψ}{\ensuremath{\Psi}}
\newunicodechar{Ω}{\ensuremath{\Omega}}
\newunicodechar{↦}{\ensuremath{\mapsto}}
\newunicodechar{∈}{\ensuremath{\in}}
\newunicodechar{∉}{\ensuremath{\notin}}
\newunicodechar{∧}{\ensuremath{\wedge}}
\newunicodechar{⇔}{\ensuremath{\Leftrightarrow}}
\newunicodechar{⟺}{\ensuremath{\Longleftrightarrow}}
\newunicodechar{⇒}{\ensuremath{\Rightarrow}}
\newunicodechar{⟹}{\ensuremath{\Longrightarrow}}
\newunicodechar{∘}{\ensuremath{\circ}}
\newunicodechar{∪}{\ensuremath{\cup}}
\newunicodechar{∩}{\ensuremath{\cap}}
\newunicodechar{≡}{\ensuremath{\equiv}}
\newunicodechar{⊂}{\ensuremath{\subset}}
\newunicodechar{⊥}{\ensuremath{\perp}}
\newunicodechar{∃}{\ensuremath{\exists}}
\newunicodechar{∀}{\ensuremath{\forall}}
\newunicodechar{∛}{\ensuremath{\sqrt[3]{\,}}}
\newunicodechar{∜}{\ensuremath{\sqrt[4]{\,}}}
\newunicodechar{⃗}{\ensuremath{{}^{\to}}}
\newunicodechar{ⁿ}{\ifmmode^{n}\else\textsuperscript{\ensuremath{n}}\fi}
\newunicodechar{ˣ}{\ifmmode^{x}\else\textsuperscript{\ensuremath{x}}\fi}
\newunicodechar{ᵇ}{\ifmmode^{b}\else\textsuperscript{\ensuremath{b}}\fi}
\newunicodechar{ʳ}{\ifmmode^{r}\else\textsuperscript{\ensuremath{r}}\fi}
\newunicodechar{ᵘ}{\ifmmode^{u}\else\textsuperscript{\ensuremath{u}}\fi}
\newunicodechar{ʸ}{\ifmmode^{y}\else\textsuperscript{\ensuremath{y}}\fi}
\newunicodechar{ᵃ}{\ifmmode^{a}\else\textsuperscript{\ensuremath{a}}\fi}
\newunicodechar{ᵉ}{\ifmmode^{e}\else\textsuperscript{\ensuremath{e}}\fi}
\newunicodechar{ᵏ}{\ifmmode^{k}\else\textsuperscript{\ensuremath{k}}\fi}
\newunicodechar{ᵖ}{\ifmmode^{p}\else\textsuperscript{\ensuremath{p}}\fi}
\newunicodechar{ᴺ}{\ifmmode^{N}\else\textsuperscript{\ensuremath{N}}\fi}
\newunicodechar{ᶜ}{\ifmmode^{c}\else\textsuperscript{\ensuremath{c}}\fi}
\newunicodechar{ʰ}{\ifmmode^{h}\else\textsuperscript{\ensuremath{h}}\fi}
\newunicodechar{ᵠ}{\ifmmode^{q}\else\textsuperscript{\ensuremath{q}}\fi}
\newunicodechar{ₙ}{\ifmmode_{n}\else\textsubscript{\ensuremath{n}}\fi}
\newunicodechar{ₐ}{\ifmmode_{a}\else\textsubscript{\ensuremath{a}}\fi}
\newunicodechar{ᵢ}{\ifmmode_{i}\else\textsubscript{\ensuremath{i}}\fi}
\newunicodechar{ᵣ}{\ifmmode_{r}\else\textsubscript{\ensuremath{r}}\fi}
\newunicodechar{ₖ}{\ifmmode_{k}\else\textsubscript{\ensuremath{k}}\fi}
\newunicodechar{ᵦ}{\ifmmode_{\beta}\else\textsubscript{\ensuremath{\beta}}\fi}
\newunicodechar{ₓ}{\ifmmode_{x}\else\textsubscript{\ensuremath{x}}\fi}
\newfontfamily\popdisplay{ArchivoBlack-Regular.ttf}[Path=../../../fonts/]
\newfontfamily\poplogo{SpaceGrotesk-Bold.ttf}[Path=../../../fonts/]
\newfontfamily\popsemi{PlusJakartaSans-SemiBold.ttf}[Path=../../../fonts/]
\newfontfamily\popxbold{PlusJakartaSans-ExtraBold.ttf}[Path=../../../fonts/]
\newfontfamily\popxboldls{PlusJakartaSans-ExtraBold.ttf}[Path=../../../fonts/,LetterSpace=3.0]

\setlist[enumerate]{leftmargin=1.7em,itemsep=5pt,topsep=4pt}
\setlist[itemize]{leftmargin=1.7em,itemsep=4pt,topsep=4pt,label={\color{poporange}\textbullet}}
\setlength{\parindent}{0pt}
\setlength{\parskip}{5pt}
\renewcommand{\arraystretch}{1.25}

% pgfplots : style Pop par défaut
\pgfplotsset{
  every axis/.append style={axis line style={popink,line width=1pt},tick style={popink},
    grid style={popline,line width=0.5pt},label style={font=\small},tick label style={font=\footnotesize}},
  popcurve/.style={poporange,line width=1.8pt,no markers,smooth},
}
% Même style disponible dans un \draw TikZ simple (hors pgfplots)
\tikzset{popcurve/.style={poporange,line width=1.8pt}}
\tikzset{popgrid/.style={popline,very thin}}
\colorlet{popcurve}{poporange} % le nom du style est parfois utilisé comme couleur

% ============================================================
% Pill / squiggle
% ============================================================
\newcommand{\poppill}[3]{%
  \tikz[baseline=-0.55ex]\node[draw=popink,line width=1.2pt,rounded corners=2.6pt,%
    fill=#2,inner xsep=6.5pt,inner ysep=3pt]{%
    \popxboldls\fontsize{7}{7}\selectfont\color{#3}\MakeUppercase{#1}};%
}
\newcommand{\popsquiggle}[1]{%
  \tikz[baseline=-0.4ex]{
    \draw[#1,line width=2.6pt,line cap=round]
      plot [smooth] coordinates {(0,0.09) (0.34,0.01) (0.68,0.21) (1.02,0.01) (1.36,0.10)};
  }%
}
% Mot-clé surligné (marqueur orange sous le texte)
\newcommand{\cle}[1]{{\popxbold\color{popdark}#1}}

% ============================================================
% En-tête / pied
% ============================================================
\pagestyle{fancy}
\fancyhf{}
\renewcommand{\headrulewidth}{0pt}
\renewcommand{\footrulewidth}{0pt}
\lhead{\raisebox{-0.15ex}{\poplogo\fontsize{13}{13}\selectfont\color{popink}OnBuch\color{poporange}+}}
\rhead{\poppill{\DOCMATIERE\ \textbullet\ \DOCNIVEAU\ \textbullet\ Leçon \DOCLECON}{white}{popink}}
\lfoot{\popxbold\fontsize{7.6}{7.6}\selectfont\color{popmuted}\MakeUppercase{Cours signé OnBuch+}\ \textbullet\ {\popsemi\DOCTITRE}}
\rfoot{\tikz[baseline=-0.6ex]\node[rounded corners=8.5pt,fill=popink,minimum height=17pt,minimum width=17pt,inner xsep=5pt,inner sep=0]{\popxbold\fontsize{8}{8}\selectfont\color{popcream}\thepage\,/\,\pageref*{LastPage}};}

% ============================================================
% Titres
% ============================================================
\newcommand{\chaptitre}[2]{%
  \begin{tcolorbox}[enhanced,colback=poporange,colframe=popink,boxrule=2.6pt,arc=11pt,
      drop shadow=popink,left=15pt,right=15pt,top=12pt,bottom=11pt,before skip=3pt,after skip=18pt]
    \poppill{#2}{popink}{popcream}\par\vspace{9pt}
    {\raggedright\hyphenpenalty=10000\exhyphenpenalty=10000\popdisplay\fontsize{22}{24}\selectfont\color{popcream}\MakeUppercase{#1}\par}
  \end{tcolorbox}
}
% Section numérotée : pastille orange avec le numéro + titre Archivo
\newcounter{popsec}
\newcommand{\coursec}[1]{%
  \stepcounter{popsec}\setcounter{popsub}{0}%
  \par\vspace{16pt}\noindent
  \begin{minipage}{\linewidth}
  \tikz[baseline=-0.7ex]\node[draw=popink,line width=1.6pt,fill=poporange,rounded corners=4pt,minimum size=22pt,inner sep=0]{\popdisplay\fontsize{12}{12}\selectfont\color{popcream}\thepopsec};%
  \hspace{8pt}{\popdisplay\fontsize{15}{17}\selectfont\color{popink}\MakeUppercase{#1}}\par
  \vspace{4pt}\noindent{\color{poporange}\rule{36pt}{3.6pt}}
  \end{minipage}\par\nopagebreak\vspace{8pt}
}
\newcounter{popsub}
\newcommand{\courssub}[1]{%
  \stepcounter{popsub}%
  \par\vspace{9pt}\noindent{\popxbold\fontsize{11.5}{13}\selectfont\color{popink}{\color{poporange}\thepopsec.\thepopsub}\hspace{6pt}#1}\par\nopagebreak\vspace{4pt}
}

% ============================================================
% Boîtes pédagogiques — même recette : fond blanc, bordure épaisse colorée,
% ombre dure encre, badge plein. Titre optionnel [#1].
% ============================================================
\tcbset{popbase/.style={breakable,enhanced,colback=white,boxrule=2.3pt,arc=9pt,
  shadow={3pt}{-3pt}{0pt}{popink},left=14pt,right=14pt,top=11pt,bottom=11pt,before skip=11pt,after skip=15pt,
  fontupper=\popsemi\fontsize{10.6}{14.5}\selectfont\color{popink},
  fontlower=\popsemi\fontsize{10.6}{14.5}\selectfont\color{popink}}}
% Coche dessinée (le glyphe ✓ n'existe pas dans les polices du document)
\newcommand{\popcheck}[1]{\tikz[baseline=-0.5ex]\draw[#1,line width=1.6pt,line cap=round,line join=round](-0.13,0.0)--(-0.04,-0.09)--(0.13,0.1);}
\providecommand{\coche}{\popcheck{popgreen}}
\providecommand{\resultat}[1]{\par\smallskip\noindent\fcolorbox{popgreen}{popgreenL}{\parbox{\dimexpr\linewidth-2\fboxsep-2\fboxrule\relax}{\textbf{Résultat.} #1}}\par\smallskip}
\newcommand{\popboxhead}[3]{\poppill{#1}{#2}{white}\ifx\relax#3\relax\else\hspace{6pt}{\popxbold\fontsize{10.6}{12}\selectfont\color{popink}#3}\fi\par\vspace{7pt}}

\newtcolorbox{aretenir}[1][]{popbase,colframe=popgreen,before upper={\popboxhead{À retenir}{popgreen}{#1}}}
% Texte en allemand (document de lecture, dialogue, extrait) : cadre
% d'étude. Un titre optionnel (source, auteur) s'affiche dans l'en-tête.
\newtcolorbox{textebox}[1][]{popbase,colframe=popblue,before upper={\popboxhead{Text}{popblue}{#1}\begin{otherlanguage}{ngerman}},after upper={\end{otherlanguage}}}
% Marques ✓ / ✗ (forme correcte / fautive) souvent utilisées par les modèles
\providecommand{\cross}{\textcolor{poppink}{\ensuremath{\times}}}
\providecommand{\xmark}{\cross}\providecommand{\crossmark}{\cross}\providecommand{\cmark}{\ensuremath{\checkmark}}
% Ligne de réponse / justification à compléter (inventée par les modèles)
\providecommand{\justification}[1]{\par\noindent\textit{Justification~:} \dotfill\par}
% \textipa (package tipa non chargé) : la police ne couvre pas l'API -> simple passage
\providecommand{\textipa}[1]{#1}
% Soulignement (ulem non chargé) : \uline, \ul
\providecommand{\uline}[1]{\underline{#1}}\providecommand{\ul}[1]{\underline{#1}}
% \glossaire{\item ...} inventée par les modèles : liste de glossaire
\providecommand{\glossaire}[1]{\begin{itemize}#1\end{itemize}}
% \alert (beamer) inventée par les modèles : texte mis en évidence
\providecommand{\alert}[1]{\textbf{\textcolor{poppink}{#1}}}
% variantes ASCII d'ordinaux écrites en macro
\providecommand{\iere}{\textsuperscript{re}}\providecommand{\ieme}{\textsuperscript{e}}\providecommand{\ere}{\textsuperscript{re}}\providecommand{\eme}{\textsuperscript{e}}\providecommand{\er}{\textsuperscript{er}}\providecommand{\ie}{\textsuperscript{e}}
% Ordinaux français écrits en macro par les modèles : 1\ière, 3\ième
\newcommand{\ière}{\textsuperscript{re}}\newcommand{\ième}{\textsuperscript{e}}
% \exemple (inventée par les modèles) : étiquette « Exemple : »
\providecommand{\exemple}{\textbf{Exemple~:}\ }
% \square utilisable aussi hors mode math (cases à cocher)
\AtBeginDocument{\let\squareMath\square\renewcommand{\square}{\ensuremath{\squareMath}}}
% Mot ou phrase en allemand à l'intérieur d'un paragraphe français
\newcommand{\all}[1]{\ifmmode\text{\textit{#1}}\else\begingroup\selectlanguage{ngerman}\itshape #1\endgroup\fi}
\let\esp\all % alias toléré
\newtcolorbox{exemplebox}[1][]{popbase,colframe=poporange,before upper={\popboxhead{Exemple}{poporange}{#1}}}
\newtcolorbox{definition}[1][]{popbase,colframe=popblue,before upper={\popboxhead{Définition}{popblue}{#1}}}
\newtcolorbox{propriete}[1][]{popbase,colframe=poppurple,before upper={\popboxhead{Propriété}{poppurple}{#1}}}
\newtcolorbox{methode}[1][]{popbase,colframe=popgold,before upper={\popboxhead{Méthode}{popgold}{#1}}}
\newtcolorbox{attention}[1][]{popbase,colframe=poppink,before upper={\popboxhead{Attention}{poppink}{#1}}}
\newtcolorbox{experience}[1][]{popbase,colframe=popblue,colback=popblueL,before upper={\popboxhead{Expérience}{popblue}{#1}}}
% « Le savais-tu ? » : carte violette pleine (contexte, culture, Cameroun)
\newtcolorbox{savaistu}[1][]{popbase,colback=poppurple,colframe=popink,
  fontupper=\popsemi\fontsize{10.4}{14.2}\selectfont\color{white},
  before upper={\poppill{Le savais-tu\,?}{white}{poppurple}\ifx\relax#1\relax\else\hspace{6pt}{\popxbold\color{white}#1}\fi\par\vspace{7pt}}}
% Exercice résolu : énoncé en haut, \tcblower puis la solution sur fond crème
\newcounter{popexo}\newcounter{popexr}
\newtcolorbox{exoresolu}[1][]{popbase,colframe=poporange,colbacklower=popcream,
  segmentation style={popink,line width=1.2pt,dashed},
  before upper={\stepcounter{popexr}\popboxhead{Exercice résolu \thepopexr}{poporange}{#1}},
  before lower={\poppill{Solution}{popink}{popcream}\par\vspace{6pt}}}
% Exercice (non résolu en place) — la correction est dans la partie Corrigés
\newtcolorbox{exercice}[1][]{popbase,colframe=popink,
  before upper={\stepcounter{popexo}\popboxhead{Exercice \thepopexo}{popink}{#1}}}
% Corrigé
\newtcolorbox{corrige}[1][]{popbase,colframe=popgreen,colback=popgreenL,
  before upper={\popboxhead{Corrigé}{popgreen}{#1}}}
% Objectifs / prérequis / bilan
\newtcolorbox{objectifs}{popbase,colframe=popink,colback=poporangeL,
  before upper={\popboxhead{Objectifs de la leçon}{popink}{}}}
\newtcolorbox{prerequis}{popbase,colframe=popink,colback=popblueL,
  before upper={\popboxhead{Prérequis}{popblue}{}}}
\newtcolorbox{bilan}{popbase,colframe=popink,colback=white,boxrule=2.8pt,
  before upper={\popboxhead{Fiche bilan}{poporange}{L'essentiel à connaître pour l'examen}}}
% Figure encadrée avec légende
\newtcolorbox{popfigure}[1][]{enhanced,breakable=false,colback=white,colframe=popline,boxrule=1.4pt,arc=8pt,
  left=10pt,right=10pt,top=10pt,bottom=8pt,before skip=10pt,after skip=12pt,halign=center,
  lower separated=false,halign lower=center,
  fontlower=\popsemi\fontsize{9}{11.5}\selectfont\color{popmuted},
  #1}
\newcommand{\legende}[1]{\par\vspace{6pt}{\popsemi\fontsize{9}{11.5}\selectfont\color{popmuted}#1}}

% ============================================================
% Signature OnBuch+
% ============================================================
\newcommand{\poplogoinline}[1]{{\poplogo\fontsize{#1}{#1}\selectfont\color{popink}OnBuch\color{poporange}+}}
\newcommand{\signatureonbuch}{%
  \begin{tcolorbox}[enhanced,colback=popink,colframe=popink,boxrule=2.2pt,arc=10pt,
      drop shadow=poporange,left=14pt,right=14pt,top=10pt,bottom=10pt,before skip=0pt,after skip=0pt]
    \tikz[baseline=-0.7ex]\node[circle,fill=popgreen,draw=popcream,line width=1.3pt,minimum size=22pt,inner sep=0]{\popcheck{white}};%
    \hspace{9pt}\begin{minipage}[c]{0.62\linewidth}
      {\popxbold\fontsize{12}{13}\selectfont\color{popcream}Signé\ \textbullet\ {\poplogo OnBuch\color{poporange}+}}\par\vspace{2pt}
      {\raggedright\popsemi\fontsize{8}{10.5}\selectfont\color{popline}\MakeUppercase{Équipe pédagogique OnBuch+}\\\MakeUppercase{Conforme au programme officiel MINESEC}\par}
    \end{minipage}\hfill
    {\poplogo\fontsize{22}{22}\selectfont\color{popcream}OnBuch\color{poporange}+}
  \end{tcolorbox}%
}

% ============================================================
% Page de garde
% #1 titre · #2 matière · #3 niveau · #4 module · #5 n° leçon · #6 durée ·
% #7 description / objectifs (texte ou itemize)
% ============================================================
\newcommand{\pagegarde}[7]{%
  \thispagestyle{empty}
  \newgeometry{a4paper,margin=1.7cm,top=1.6cm,bottom=1.4cm}%
  \begin{tikzpicture}[remember picture,overlay]
    \fill[poporange!18] ([xshift=-3.2cm,yshift=-3.6cm]current page.north east) circle (4.6cm);
    \fill[poppurple!14] ([xshift=2.4cm,yshift=7.6cm]current page.south west) circle (3.8cm);
  \end{tikzpicture}%
  {\poplogo\fontsize{30}{30}\selectfont\color{popink}OnBuch\color{poporange}+}\hfill
  \raisebox{0.6ex}{\poppill{Cours signé \textbullet\ Premium}{popink}{popcream}}\par
  \vspace{1pt}\popsquiggle{poppurple}\par
  \vspace{8pt}
  \begin{tcolorbox}[enhanced,colback=poporange,colframe=popink,boxrule=2.8pt,arc=13pt,
      drop shadow=popink,left=18pt,right=18pt,top=12pt,bottom=13pt,after skip=10pt]
    \poppill{#2 \textbullet\ #3}{popink}{popcream}\hspace{5pt}\poppill{#4}{white}{popink}\par\vspace{8pt}
    {\popdisplay\fontsize{14}{14}\selectfont\color{popink}LEÇON #5}\par\vspace{4pt}
    {\raggedright\hyphenpenalty=10000\exhyphenpenalty=10000\popdisplay\fontsize{25}{27}\selectfont\color{popcream}\MakeUppercase{#1}\par}
    \vspace{8pt}
    {\popxbold\fontsize{9.5}{9.5}\selectfont\color{popink}\MakeUppercase{Durée conseillée : #6}}
  \end{tcolorbox}
  \begin{tcolorbox}[enhanced,colback=white,colframe=popink,boxrule=2.2pt,arc=10pt,
      drop shadow=popink,left=16pt,right=16pt,top=10pt,bottom=10pt,after skip=8pt]
    \poppill{Dans cette leçon}{poppurple}{white}\par\vspace{5pt}
    \popsemi\fontsize{9.3}{12.6}\selectfont\color{popink}#7
  \end{tcolorbox}
  \noindent\begin{minipage}[t]{0.487\linewidth}
    \begin{tcolorbox}[enhanced,colback=popgreen,colframe=popink,boxrule=2.2pt,arc=10pt,
        drop shadow=popink,left=11pt,right=11pt,top=10pt,bottom=10pt,height=3.1cm,valign=center]
      \tcbox[colback=white,colframe=popink,boxrule=1.3pt,arc=5pt,boxsep=0pt,nobeforeafter,
        left=3pt,right=3pt,top=3pt,bottom=3pt,valign=center]{\qrcode[height=1.9cm]{https://play.google.com/store/apps/details?id=com.luvvix.onbuch&hl=fr}}%
      \hspace{8pt}\begin{minipage}[c]{0.5\linewidth}
        {\popdisplay\fontsize{11}{12.5}\selectfont\color{white}\MakeUppercase{Télécharge OnBuch}}\par\vspace{4pt}
        {\raggedright\popsemi\fontsize{8.4}{11}\selectfont\color{white}Gratuit \textbullet\ Google Play\\Tuteur IA, annales, concours, résultats\par}
      \end{minipage}
    \end{tcolorbox}
  \end{minipage}\hfill
  \begin{minipage}[t]{0.487\linewidth}
    \begin{tcolorbox}[enhanced,colback=poppurple,colframe=popink,boxrule=2.2pt,arc=10pt,
        drop shadow=popink,left=11pt,right=11pt,top=10pt,bottom=10pt,height=3.1cm,valign=center]
      \tikz[baseline=-0.7ex]\node[circle,fill=white,draw=popink,line width=1.3pt,minimum size=1.6cm,inner sep=0]{\popdisplay\fontsize{24}{24}\selectfont\color{poppurple}+};%
      \hspace{8pt}\begin{minipage}[c]{0.6\linewidth}
        {\popdisplay\fontsize{11}{12.5}\selectfont\color{white}\MakeUppercase{Passe à OnBuch+}}\par\vspace{4pt}
        {\raggedright\popsemi\fontsize{8.4}{11}\selectfont\color{white}Tous les cours signés, Tuteur IA illimité, fiches PDF — onglet OnBuch+ de l'app\par}
      \end{minipage}
    \end{tcolorbox}
  \end{minipage}
  \vspace{8pt}
  \popcontenu
  \vfill
  \signatureonbuch
  \restoregeometry
  \newpage
}

% Rangée « Ce cours contient » (page de garde)
\newcommand{\poptile}[3]{% #1 couleur #2 gros texte #3 légende
  \begin{tcolorbox}[enhanced,colback=#1,colframe=popink,boxrule=2pt,arc=9pt,shadow={2.5pt}{-2.5pt}{0pt}{popink},
      left=6pt,right=6pt,top=9pt,bottom=9pt,halign=center,valign=center,height=1.75cm,width=\linewidth,nobeforeafter]
    {\popdisplay\fontsize{12.5}{13}\selectfont\color{white}\MakeUppercase{#2}}\par\vspace{3pt}
    {\centering\popsemi\fontsize{7.8}{9.5}\selectfont\color{white}#3\par}
  \end{tcolorbox}}
\newcommand{\popcontenu}{%
  \par\noindent{\popxboldls\fontsize{7.5}{7.5}\selectfont\color{popmuted}CE COURS SIGNÉ CONTIENT}\par\vspace{7pt}
  \noindent\begin{minipage}[t]{0.235\linewidth}\poptile{popblue}{Cours}{complet, illustré, pas à pas}\end{minipage}\hfill
  \begin{minipage}[t]{0.235\linewidth}\poptile{poporange}{Méthodes}{et exercices résolus}\end{minipage}\hfill
  \begin{minipage}[t]{0.235\linewidth}\poptile{poppink}{Exercices}{type examen + corrigés}\end{minipage}\hfill
  \begin{minipage}[t]{0.235\linewidth}\poptile{popgreen}{Fiche bilan}{l'essentiel à retenir}\end{minipage}\par}

% Sommaire
\newtcolorbox{sommaire}{enhanced,colback=white,colframe=popink,boxrule=2.2pt,arc=10pt,
  drop shadow=popink,left=18pt,right=18pt,top=13pt,bottom=13pt,before skip=6pt,after skip=20pt,
  before upper={\poppill{Sommaire}{poppurple}{white}\par\vspace{9pt}\popsemi\fontsize{10.6}{14.5}\selectfont\color{popink}}}

% Signature de fin de cours — poussée tout en bas de la dernière page
\newcommand{\findecours}{%
  \par\vfill\nopagebreak
  \begin{center}\popsquiggle{poporange}\end{center}\vspace{4pt}
  \signatureonbuch
}
\AtBeginDocument{\def\llbracket{\mathopen{[\mkern-3.5mu[}}\def\rrbracket{\mathclose{]\mkern-3.5mu]}}} % intervalles d'entiers (glyphe ⟦ absent de la police)
% Glyphes absents de Fira Math : redéfinis à partir de glyphes disponibles
\AtBeginDocument{%
  \def\setminus{\mathbin{\backslash}}\let\smallsetminus\setminus
  \def\vdots{\mathord{\vbox{\baselineskip3.2pt\lineskiplimit0pt\kern4pt\hbox{.}\hbox{.}\hbox{.}}}}%
  \def\ddots{\mathinner{\mkern1mu\raise6.5pt\hbox{.}\mkern2mu\raise3.6pt\hbox{.}\mkern2mu\raise0.7pt\hbox{.}\mkern1mu}}%
  \def\triangle{\mathord{\scalebox{0.9}{$\symup{\Delta}$}}}%
  \def\nabla{\mathord{\raisebox{\depth}{\scalebox{1}[-1]{$\symup{\Delta}$}}}}%
  \def\checkmark{\popcheck{popdark}}%
  \def\star{\mathbin{\ast}}\def\bigstar{\mathord{\ast}}%
  \def\diamond{\mathbin{\scalebox{0.6}{$\Diamond$}}}%
  \def\bigcup{\mathop{\mathchoice{\raisebox{-0.3em}{\scalebox{1.7}{$\cup$}}}{\scalebox{1.25}{$\cup$}}{\cup}{\cup}}\displaylimits}%
  \def\bigcap{\mathop{\mathchoice{\raisebox{-0.3em}{\scalebox{1.7}{$\cap$}}}{\scalebox{1.25}{$\cap$}}{\cap}{\cap}}\displaylimits}%
  \def\lesseqqgtr{\mathrel{\lessgtr}}\def\bigoplus{\mathop{\oplus}\displaylimits}%
}
\providecommand{\ssi}{si et seulement si} % abréviation fréquente
\AtBeginDocument{\providecommand{\wideparen}{\overparen}} % arc de cercle

%% FIGFIX-BEGIN v5 — correctifs globaux des figures (docs/onbuch-generation/tools/figfix)
\usepackage{adjustbox}\usepackage{etoolbox}\tcbuselibrary{raster}
% toute figure TikZ/pgfplots plus large que la ligne est réduite à la largeur disponible
\BeforeBeginEnvironment{tikzpicture}{\begin{adjustbox}{max width=\linewidth}}
\AfterEndEnvironment{tikzpicture}{\end{adjustbox}}
% légendes pgfplots lisibles par-dessus les courbes
\pgfplotsset{every axis/.append style={legend style={fill=white,fill opacity=0.92,text opacity=1,draw=black!15,inner sep=3pt}}}
% étiquettes d'arêtes (syntaxe edge["..."]) avec fond blanc
\tikzset{every edge quotes/.append style={fill=white,inner sep=1.2pt}}
%% FIGFIX-END
\begin{document}

\pagegarde{Les fléaux sociaux et leurs conséquences}{Allemand}{1ère A}{Chapitre 1 : Vie familiale et sociale}{ALA03}{2 h}{À partir d'un texte sur les fléaux sociaux (pauvreté, drogue, violence, chômage), les élèves enrichissent leur lexique thématique, découvrent l'expression de la cause avec weil et denn, révisent le comparatif et le parfait, puis produisent un court texte argumenté sur les conséquences de ces problèmes dans leur environnement.
\vspace{4pt}
\begin{multicols}{2}\small
\begin{itemize}
\item À la fin de cette leçon, je sais comprendre globalement et en détail un texte allemand sur les fléaux sociaux et leurs conséquences.
\item À la fin de cette leçon, je sais employer correctement le vocabulaire du thème (die Armut, die Drogen, die Gewalt, die Arbeitslosigkeit, die Folge, gefährlich) avec articles et pluriels.
\item À la fin de cette leçon, je sais exprimer la cause avec weil (subordonnée) et denn (coordination) en plaçant le verbe correctement.
\item À la fin de cette leçon, je sais former et employer le comparatif des adjectifs (gefährlich → gefährlicher, gut → besser).
\item À la fin de cette leçon, je sais conjuguer un verbe au parfait (haben/sein + Partizip II) pour raconter des faits passés.
\item À la fin de cette leçon, je sais répondre par écrit et oralement à des questions de compréhension sur un texte.
\end{itemize}\end{multicols}}

\begin{sommaire}
\begin{multicols}{2}
\begin{enumerate}
\item Texte support : Soziale Probleme in unserer Welt
\item Compréhension du texte
\item Lexique thématique : les fléaux sociaux
\item Grammaire 1 : la cause avec weil et denn
\item Grammaire 2 : le comparatif — Conjugaison : le parfait
\item Méthode du commentaire et production guidée
\item Activité d'intégration
\item Méthodes et exercices résolus
\item Exercices
\item Corrigés des exercices
\item Fiche bilan
\end{enumerate}
\end{multicols}
\end{sommaire}

\begin{objectifs}
\begin{itemize}
\item À la fin de cette leçon, je sais comprendre globalement et en détail un texte allemand sur les fléaux sociaux et leurs conséquences.
\item À la fin de cette leçon, je sais employer correctement le vocabulaire du thème (die Armut, die Drogen, die Gewalt, die Arbeitslosigkeit, die Folge, gefährlich) avec articles et pluriels.
\item À la fin de cette leçon, je sais exprimer la cause avec weil (subordonnée) et denn (coordination) en plaçant le verbe correctement.
\item À la fin de cette leçon, je sais former et employer le comparatif des adjectifs (gefährlich → gefährlicher, gut → besser).
\item À la fin de cette leçon, je sais conjuguer un verbe au parfait (haben/sein + Partizip II) pour raconter des faits passés.
\item À la fin de cette leçon, je sais répondre par écrit et oralement à des questions de compréhension sur un texte.
\item À la fin de cette leçon, je sais résumer un texte en quelques phrases avec mes propres mots.
\item À la fin de cette leçon, je sais rédiger un court paragraphe guidé sur un fléau social et ses conséquences.
\end{itemize}
\end{objectifs}

\begin{prerequis}
\begin{itemize}
\item Le présent de l'indicatif des verbes forts et faibles (Seconde)
\item La place du verbe en position 2 dans la phrase déclarative et l'inversion du sujet
\item Les articles définis et indéfinis au nominatif et à l'accusatif
\item Le vocabulaire de base de la famille et de la société (die Familie, die Schule, die Arbeit)
\item Notions élémentaires sur le Partizip II des verbes faibles (début d'année de Première)
\end{itemize}
\end{prerequis}

\newpage

\coursec{Texte support : Soziale Probleme in unserer Welt}

\courssub{Situation d'accroche et problématique}
À Yaoundé comme à Berlin, les journaux parlent chaque jour de chômage, de drogue et de violence. En Allemagne, des associations comme \all{die Drogenhilfe} aident les jeunes en difficulté, tandis qu'au Cameroun, des ONG luttent contre la pauvreté dans les quartiers défavorisés. Comment parler de ces problèmes en allemand ? Comment expliquer pourquoi ils existent et quelles sont leurs conséquences ? Cette leçon nous donne les mots et les outils grammaticaux pour en débattre.

\courssub{Première approche du texte}
\begin{methode}[Comment lire un texte allemand]
\begin{enumerate}
\item Lire le titre et regarder l'illustration.
\item Faire une première lecture globale sans dictionnaire.
\item Repérer les mots transparents (cognats) et les mots déjà connus.
\item Faire une deuxième lecture détaillée en s'aidant du glossaire.
\item Vérifier sa compréhension phrase par phrase (sujet – verbe – compléments).
\end{enumerate}
\end{methode}

\begin{textebox}[Soziale Probleme in unserer Welt]
In vielen Ländern gibt es große soziale Probleme. Viele Menschen leben in Armut, weil sie keine Arbeit haben. Die Arbeitslosigkeit ist heute größer als früher. Manche Jugendliche nehmen Drogen, denn sie haben keine Perspektive. Drogen sind gefährlicher als Alkohol: Sie machen süchtig und krank. Auch die Gewalt in den Städten hat zugenommen. Letztes Jahr haben viele junge Menschen die Schule verlassen, weil ihre Familien zu arm waren. Die Folgen sind schwer: mehr Kriminalität, mehr Krankheiten, mehr Obdachlosigkeit.
\end{textebox}

\courssub{Glossaire bilingue (mots difficiles)}
\begin{center}
\begin{tabularx}{\linewidth}{|X|X|X|}
\hline
\rowcolor{popblueL}
\textbf{Deutsch} & \textbf{Français} & \textbf{Remarque} \\
\hline
die Armut, die Ärmlichkeiten & la misère, la pauvreté & nom féminin, pluriel en \emph{-en} \\
die Arbeitslosigkeit & le chômage & nom féminin, pas de pluriel courant \\
die Drogen (Pl.) & les drogues & toujours au pluriel \\
die Gewalt & la violence & nom féminin, pas de pluriel \\
die Folge, die Folgen & la conséquence & nom féminin \\
gefährlich & dangereux & adjectif \\
süchtig & dépendant (à une drogue) & adjectif \\
die Obdachlosigkeit & le sans‑abrisme & nom féminin \\
die Perspektive, die Perspektiven & la perspective & nom féminin \\
die Kriminalität & la criminalité & nom féminin \\
\hline
\end{tabularx}
\end{center}

\courssub{Repérage des mots transparents et hypothèses}
\begin{itemize}
\item \all{soziale Probleme} → problèmes sociaux (transparent)
\item \all{Armut} → misère (proche de \emph{arm} = pauvre)
\item \all{Arbeitslosigkeit} → chômage (Arbeit + los + igkeit)
\item \all{Drogen} → drogues (cognat)
\item \all{Gewalt} → violence (pas transparent, à mémoriser)
\item \all{Folgen} → conséquences (proche de \emph{suivre})
\end{itemize}
À partir du titre \emph{Soziale Probleme in unserer Welt} et de la carte mentale ci‑dessous, on suppose que l'article traite de la pauvreté, du chômage, de la drogue et de la violence, ainsi que de leurs conséquences sur les jeunes.

\courssub{Carte mentale du vocabulaire thématique}
\begin{popfigure}
\begin{tikzpicture}[
  mindmap,
  concept color=popblue,
  level 1/.style={level distance=3.5cm,sibling angle=90},
  level 2/.style={level distance=2.5cm,sibling angle=45},
  every node/.style={font=\small},
  text=white
]
\node[concept] {Soziale Probleme}
  child[concept color=poporange] { node[concept] {die Armut}
    child[concept color=poporange!70] { node[concept] {die Folge: Elend} } }
  child[concept color=popgreen] { node[concept] {die Arbeitslosigkeit}
    child[concept color=popgreen!70] { node[concept] {die Folge: Perspektivlosigkeit} } }
  child[concept color=poppurple] { node[concept] {die Drogen}
    child[concept color=poppurple!70] { node[concept] {die Folge: Sucht \& Krankheit} } }
  child[concept color=poppink] { node[concept] {die Gewalt}
    child[concept color=poppink!70] { node[concept] {die Folge: Kriminalität} } };
\end{tikzpicture}
\legende{Carte mentale : les quatre fléaux sociaux principaux et une conséquence typique pour chacun.}
\end{popfigure}

\courssub{Lecture à voix haute (compréhension orale)}
L'enseignant lit le texte lentement, en marquant les liaisons et l'intonation des phrases subordonnées (\all{weil}, \all{denn}). Les élèves suivent sur leur copie et notent les mots qu'ils reconnaissent à l'oral.

\courssub{Exercice résolu : compréhension globale}
\begin{exoresolu}[Question 1]
\textbf{Énoncé :} D'après le texte, citez deux causes de la pauvreté mentionnées explicitement.
\tcblower
\textbf{Solution détaillée :}
\begin{enumerate}
\item \all{weil sie keine Arbeit haben} → parce qu'ils n'ont pas de travail (chômage).
\item \all{weil ihre Familien zu arm waren} → parce que leurs familles étaient trop pauvres (pauvreté familiale).
\end{enumerate}
Les deux subordonnées causatives introduites par \all{weil} donnent les raisons directes.
\end{exoresolu}

\courssub{Attention : piège de la traduction mot à mot}
\begin{attention}[Ne pas traduire mot à mot !]
En allemand, l'ordre des mots dans la subordonnée \all{weil} place le verbe à la fin : \all{weil sie keine Arbeit haben}. En français, on dit « parce qu'ils n'ont pas de travail » (verbe avant le complément). Traduire littéralement « parce qu'ils aucun travail ont » est une erreur fréquente. Cherchez d'abord le sens global de la phrase (sujet – verbe – compléments) avant de transposer.
\end{attention}

\courssub{Le saviez‑vous ?}
\begin{savaistu}[Pauvreté en Allemagne]
En Allemagne, environ 16 \% de la population est menacée par la pauvreté (\all{armutsgefährdet}). Le mot \all{die Armut} dérive de l'adjectif \all{arm} (pauvre) + suffixe \emph{-ut} (nominalisation). En Suisse, on parle aussi de \all{Armut} mais le seuil de pauvreté y est plus élevé qu'en Allemagne.
\end{savaistu}

\coursec{Compréhension du texte}

Dans la section précédente, nous avons découvert le texte \all{Soziale Probleme in unserer Welt}, un article de journal qui décrit les grands fléaux sociaux — la pauvreté, le chômage, la drogue, la violence — et leurs conséquences. Nous allons maintenant le lire de manière méthodique : d'abord une compréhension globale, puis une compréhension détaillée, enfin un résumé guidé. C'est exactement la démarche exigée à l'examen.

\courssub{Le texte support (rappel)}

\begin{textebox}[Soziale Probleme in unserer Welt — Artikel aus der Zeitung „Der Blick“]
Viele Menschen in der Welt leben in Armut. Sie haben kein Geld für Essen, Kleidung oder eine Wohnung. Oft finden sie keine Arbeit, weil sie keine Ausbildung haben. Die Arbeitslosigkeit ist heute größer als früher, denn viele Fabriken schließen.

Auch in Kamerun sehen wir diese Probleme: In Yaoundé und Douala suchen viele Jugendliche eine Arbeit, aber sie finden keine. Manche Jugendliche nehmen Drogen, weil sie keine Perspektive haben. Drogen machen süchtig und krank. Sie zerstören Familien und Freundschaften.

Die Gewalt in den Städten ist auch gefährlich. Sie hat viele Ursachen: Armut, Arbeitslosigkeit und Drogen. Die Folgen der sozialen Probleme sind schlimm: Kinder gehen nicht zur Schule, Familien haben kein Geld, und viele Menschen sind krank.

Aber es gibt Hoffnung: Schulen, Vereine und Familien können helfen. Bildung ist der beste Weg aus der Armut.
\end{textebox}

\textbf{Glossaire (Wörterbuch) :}

\begin{center}
\begin{tabularx}{\linewidth}{|X|X|}\hline
\rowcolor{popblueL} \textbf{Deutsch} & \textbf{Français} \\ \hline
die Armut (nur Singular) & la pauvreté \\ \hline
die Arbeitslosigkeit (nur Singular) & le chômage \\ \hline
die Droge, -n & la drogue \\ \hline
die Gewalt (nur Singular) & la violence \\ \hline
die Perspektive, -n & la perspective, l'avenir \\ \hline
die Folge, -n & la conséquence \\ \hline
die Ursache, -n & la cause \\ \hline
süchtig machen & rendre dépendant \\ \hline
zerstören & détruire \\ \hline
gefährlich & dangereux \\ \hline
die Hoffnung, -en & l'espoir \\ \hline
der Weg, -e & le chemin, le moyen \\ \hline
\end{tabularx}
\end{center}

\begin{savaistu}[Le Cameroun et l'Allemagne face aux défis sociaux]
Le texte mentionne Yaoundé et Douala. En Allemagne, on parle de \all{Armut trotz Arbeit} (pauvreté malgré le travail) et de \all{Jugendarbeitslosigkeit} (chômage des jeunes). La \all{Bundesagentur für Arbeit} accompagne les demandeurs d'emploi. Au Cameroun, le \all{Programme d'Appui à l'Insertion Professionnelle des Jeunes} (PAIPJ) et les \all{Écoles de Deuxième Chance} tentent d'offrir des perspectives. L'éducation (\all{Bildung}) est partout la clé : \all{Bildung ist der beste Weg aus der Armut}.
\end{savaistu}

\courssub{La démarche de compréhension}

Avant de répondre aux questions, voici la méthode complète, étape par étape.

\begin{popfigure}
\begin{tikzpicture}[
  node distance=0.55cm,
  box/.style={draw=popblue, fill=popblueL, rounded corners, align=center, text width=4.6cm, minimum height=0.9cm, font=\small},
  arr/.style={-{Stealth[length=2.5mm]}, thick, poporange}
]
\node[box, fill=popgoldL, draw=popgold, align=center] (t){\textbf{1. Titre et source}\\De quoi va parler le texte ?};
\node[box, below=of t, align=center] (lg){\textbf{2. Lecture globale}\\Lire vite, sans dictionnaire};
\node[box, below=of lg, align=center] (qg){\textbf{3. Questions globales}\\Type de texte ? Thème ? Qui ?};
\node[box, below=of qg, align=center] (ld){\textbf{4. Lecture détaillée}\\Souligner les mots-clés};
\node[box, below=of ld, align=center] (qd){\textbf{5. Questions détaillées}\\Richtig / Falsch + citations};
\node[box, below=of qd, fill=popgreenL, draw=popgreen, align=center] (r){\textbf{6. Résumé}\\Anfang – Mitte – Schluss};
\draw[arr] (t) -- (lg);
\draw[arr] (lg) -- (qg);
\draw[arr] (qg) -- (ld);
\draw[arr] (ld) -- (qd);
\draw[arr] (qd) -- (r);
\end{tikzpicture}
\legende{Organigramme de la démarche de compréhension d'un texte}
\end{popfigure}

\courssub{Compréhension globale}

Après une première lecture rapide, on répond aux questions générales :

\begin{itemize}
\item \textbf{Worum geht es im Text?} (De quoi parle le texte ?) — \all{Der Text handelt von sozialen Problemen und ihren Folgen.} (Le texte traite des problèmes sociaux et de leurs conséquences.)
\item \textbf{Welcher Texttyp ist das?} (Quel est le type de texte ?) — \all{Das ist ein Zeitungsartikel.} (C'est un article de journal.) On le voit à la source indiquée : \all{aus der Zeitung „Der Blick“}.
\item \textbf{Wer ist betroffen?} (Qui est concerné ?) — \all{Viele Menschen in der Welt, besonders Jugendliche in den Städten, auch in Kamerun.} (Beaucoup de personnes dans le monde, surtout les jeunes des villes, y compris au Cameroun.)
\end{itemize}

\begin{methode}[Répondre à une question de compréhension]
\begin{enumerate}
\item \textbf{Repérer} la ligne du texte où se trouve l'information (souligner le passage).
\item \textbf{Recopier} mentalement l'information essentielle.
\item \textbf{Reformuler} en une phrase complète, avec le verbe conjugué en \cle{deuxième position}.
\item \textbf{Vérifier} : la phrase commence par une majuscule, se termine par un point, et tous les noms prennent une majuscule.
\end{enumerate}
\end{methode}

\begin{attention}[Jamais une réponse en un seul mot !]
En allemand, on répond toujours par une \cle{phrase complète}, jamais par un mot isolé. Si la question est \all{Warum nehmen Jugendliche Drogen?}, la réponse \all{Keine Perspektive.} est \textbf{fausse} à l'examen. Il faut écrire : \all{Sie nehmen Drogen, weil sie keine Perspektive haben.} Et attention : le verbe conjugué est en deuxième position (\all{Sie \textbf{nehmen} Drogen}), et après \all{weil}, il va à la fin (\all{weil sie keine Perspektive \textbf{haben}}).
\end{attention}

\courssub{Compréhension détaillée : les questions}

\begin{exemplebox}[Questions et réponses modèles]
\all{\textbf{Warum leben viele Menschen in Armut?}} (Pourquoi beaucoup de gens vivent-ils dans la pauvreté ?)\par
→ \all{Viele Menschen leben in Armut, weil sie keine Arbeit und keine Ausbildung haben.} « Beaucoup de gens vivent dans la pauvreté parce qu'ils n'ont ni travail ni formation. »\par\medskip
\all{\textbf{Was ist heute größer als früher?}} (Qu'est-ce qui est plus grand aujourd'hui qu'autrefois ?)\par
→ \all{Die Arbeitslosigkeit ist heute größer als früher.} « Le chômage est plus important aujourd'hui qu'autrefois. »\par\medskip
\all{\textbf{Warum nehmen manche Jugendliche Drogen?}} (Pourquoi certains jeunes prennent-ils de la drogue ?)\par
→ \all{Sie nehmen Drogen, weil sie keine Perspektive haben.} « Ils prennent de la drogue parce qu'ils n'ont pas de perspective. »\par\medskip
\all{\textbf{Was sind die Folgen der sozialen Probleme?}} (Quelles sont les conséquences des problèmes sociaux ?)\par
→ \all{Kinder gehen nicht zur Schule, Familien haben kein Geld, und viele Menschen sind krank.} « Les enfants ne vont pas à l'école, les familles n'ont pas d'argent et beaucoup de gens sont malades. »
\end{exemplebox}

\courssub{Richtig oder Falsch ? Justifier par une citation}

Pour les questions vrai/faux, il faut \textbf{toujours} citer la phrase du texte qui prouve la réponse.

\begin{exoresolu}[Richtig oder Falsch ?]
Énoncé : réponds par \all{Richtig} ou \all{Falsch} et cite le texte.\par
1. \all{Alle Jugendlichen in Kamerun nehmen Drogen.}\par
2. \all{Die Gewalt hat nur eine Ursache.}\par
3. \all{Bildung kann aus der Armut helfen.}
\tcblower
Solution détaillée :\par
1. \all{\textbf{Falsch.} Der Text sagt: „\textbf{Manche} Jugendliche nehmen Drogen.“} — Le texte dit « certains jeunes », pas « tous ». Piège classique : \all{manche} (certains) n'est pas \all{alle} (tous).\par
2. \all{\textbf{Falsch.} Der Text sagt: „Sie hat \textbf{viele} Ursachen: Armut, Arbeitslosigkeit und Drogen.“} — La violence a plusieurs causes, pas une seule.\par
3. \all{\textbf{Richtig.} Der Text sagt: „Bildung ist der beste Weg aus der Armut.“} — Citation exacte : l'éducation est le meilleur moyen de sortir de la pauvreté.
\end{exoresolu}

\courssub{Recherche dans le texte : les champs lexicaux}

On demande souvent de retrouver les mots d'un champ lexical. Voici le résultat pour notre texte :

\begin{center}
\begin{tabularx}{\linewidth}{|X|X|}\hline
\rowcolor{popblueL} \textbf{Champ lexical des problèmes sociaux} & \textbf{Champ lexical des conséquences} \\ \hline
\all{die Armut} (la pauvreté) & \all{süchtig und krank} (dépendant et malade) \\ \hline
\all{die Arbeitslosigkeit} (le chômage) & \all{zerstören Familien} (détruisent les familles) \\ \hline
\all{die Drogen} (les drogues) & \all{Kinder gehen nicht zur Schule} (les enfants ne vont pas à l'école) \\ \hline
\all{die Gewalt} (la violence) & \all{Familien haben kein Geld} (les familles n'ont pas d'argent) \\ \hline
\all{keine Perspektive} (pas d'avenir) & \all{viele Menschen sind krank} (beaucoup de gens sont malades) \\ \hline
\end{tabularx}
\end{center}

\begin{aretenir}[Mots-clés pour l'examen]
Retiens ces paires \cle{nom + verbe} ou \cle{nom + adjectif} fréquentes dans les sujets de bac : \all{Armut bekämpfen}, \all{Arbeitslosigkeit steigen}, \all{Drogen nehmen / süchtig werden}, \all{Gewalt ausüben / verhindern}, \all{Folgen haben / tragen}, \all{Ursachen suchen / nennen}, \all{Perspektive schaffen / fehlen}, \all{Hoffnung geben / verlieren}.
\end{aretenir}

\courssub{Reformulation : expliquer une phrase clé}

Expliquer avec ses propres mots, c'est montrer qu'on a \emph{compris}, pas seulement recopié. Prenons la phrase clé : \all{„Drogen machen süchtig und krank.“}

\begin{methode}[Comment reformuler une phrase]
\begin{enumerate}
\item Identifier le(s) mot(s)-clé(s) difficiles (\all{süchtig}, \all{krank}).
\item Les remplacer par une explication ou un synonyme simple (\all{nicht aufhören können}, \all{Gesundheit kaputt}).
\item Reconstruire une phrase grammaticalement correcte (sujet + verbe 2ème position).
\end{enumerate}
\end{methode}

\begin{exemplebox}[Reformulations possibles]
\all{Wenn man Drogen nimmt, kann man nicht mehr aufhören, und der Körper wird krank.} « Quand on prend de la drogue, on ne peut plus s'arrêter, et le corps tombe malade. »\par\medskip
\all{Drogen sind gefährlich: Man braucht sie immer wieder, und die Gesundheit ist kaputt.} « Les drogues sont dangereuses : on en a toujours besoin à nouveau, et la santé est détruite. »
\end{exemplebox}

\courssub{Résumé guidé : Anfang – Mitte – Schluss}

\begin{methode}[Faire un résumé en trois phrases]
\begin{enumerate}
\item Garder \textbf{uniquement les idées principales} (une pour le début, une pour le milieu, une pour la fin).
\item Utiliser des \cle{connecteurs} : \all{zuerst} (d'abord), \all{dann} (ensuite), \all{schließlich} (finalement).
\item \textbf{Ne jamais recopier} le texte mot à mot : changer les formulations.
\item Rester court : trois phrases maximum.
\end{enumerate}
\end{methode}

\begin{exemplebox}[Modèle de résumé]
\all{Zuerst zeigt der Text, dass viele Menschen arm und arbeitslos sind. Dann erklärt er, dass manche Jugendliche Drogen nehmen und dass Gewalt gefährlich ist. Schließlich sagt er, dass Bildung die beste Hilfe ist.}\par
« D'abord, le texte montre que beaucoup de gens sont pauvres et au chômage. Ensuite, il explique que certains jeunes prennent de la drogue et que la violence est dangereuse. Finalement, il dit que l'éducation est la meilleure aide. »
\end{exemplebox}

\begin{attention}[Erreurs fréquentes des francophones]
\begin{itemize}
\item Ne traduis pas « d'abord » par \all{am ersten} : dis \all{zuerst} (ou \all{als erstes}).
\item Attention à la place du verbe avec les connecteurs : \all{zuerst} et \all{dann} comptent comme premier élément, donc le verbe vient juste après : \all{\textbf{Dann erklärt er}...} (et non \all{Dann er erklärt}).
\item \all{schließlich} ne signifie pas « schließlich = finalement » au sens de « enfin ! » (soulagement) ; c'est bien « finalement » au sens de « pour conclure ».
\item Après \all{dass} (conjonction de subordination), le verbe va à la \cle{fin de la phrase} : \all{..., dass viele Menschen \textbf{sind}.}
\end{itemize}
\end{attention}

\begin{exercice}[À toi de jouer !]
1. Réponds par des phrases complètes : \all{Warum ist die Arbeitslosigkeit heute größer? Wo suchen viele Jugendliche eine Arbeit?}\par
2. \all{Richtig oder Falsch?} Justifie par une citation : \all{In Douala finden alle Jugendlichen eine Arbeit.}\par
3. Résume le texte en trois phrases avec \all{zuerst}, \all{dann}, \all{schließlich}.
\end{exercice}

\begin{corrige}[Exercice « À toi de jouer ! »]
1. \all{Die Arbeitslosigkeit ist heute größer, denn viele Fabriken schließen. Viele Jugendliche suchen in Yaoundé und Douala eine Arbeit.}\par
2. \all{\textbf{Falsch.} Der Text sagt: „\textbf{Sie suchen eine Arbeit, aber sie finden keine.}“}\par
3. Exemple : \all{Zuerst spricht der Text über Armut und Arbeitslosigkeit. Dann zeigt er die Probleme der Jugendlichen: Drogen und Gewalt. Schließlich gibt er Hoffnung: Bildung hilft.}
\end{corrige}

\coursec{Lexique thématique : les fléaux sociaux}

Dans la section précédente, nous avons lu et compris le texte \all{Soziale Probleme in unserer Welt}. Nous y avons rencontré des mots comme \all{die Armut}, \all{die Arbeitslosigkeit} ou \all{die Drogen}. Pour pouvoir parler et écrire correctement sur les fléaux sociaux, il faut maintenant maîtriser ce vocabulaire de façon systématique : article, pluriel, famille de mots et emploi dans une phrase. C'est l'objectif de cette section.

\begin{methode}[Comment apprendre un nom allemand]
Pour chaque nom, retiens toujours trois informations ensemble :
\begin{enumerate}
\item \cle{le mot} lui-même : \all{Armut} ;
\item \cle{l'article} (le genre) : \all{die} ;
\item \cle{le pluriel} (s'il existe) : ici aucun, on note (–).
\end{enumerate}
On écrit donc : \all{die Armut (–)}. Apprends toujours le mot dans une phrase complète : \all{Viele Menschen leben in Armut.} « Beaucoup de gens vivent dans la pauvreté. »
\end{methode}

\courssub{Champ 1 : les problèmes sociaux}

\begin{definition}[die Arbeitslosigkeit]
\all{Die Arbeitslosigkeit (–)} signifie « le chômage ». C'est un mot composé : \all{arbeits-} (travail) + \all{-los} (sans) + \all{-igkeit} (suffixe formant un nom féminin abstrait). Littéralement : « l'état d'être sans travail ». Comme tous les noms en \all{-igkeit}, il est féminin et n'a pas de pluriel.
\end{definition}

\begin{center}
\begin{tabularx}{\linewidth}{|l|X|X|}
\hline
\rowcolor{popblueL} Deutsch & Pluriel & Français \\
\hline
die Armut & – & la pauvreté \\
\hline
die Arbeitslosigkeit & – & le chômage \\
\hline
die Drogen & (presque toujours au pluriel) & la drogue, les drogues \\
\hline
die Gewalt & – & la violence \\
\hline
die Kriminalität & – & la criminalité \\
\hline
die Obdachlosigkeit & – & le sans-abrisme \\
\hline
der Hunger & – & la faim \\
\hline
\end{tabularx}
\end{center}

\begin{attention}[die Drogen : presque toujours au pluriel]
En allemand, \all{die Drogen} s'emploie presque toujours au pluriel quand on parle du problème en général : \all{Drogen sind gefährlich.} « La drogue est dangereuse. » Au singulier, \all{die Droge (–n)} désigne une substance précise. Ne confonds pas avec \all{das Medikament (–e)} « le médicament », qui est un faux ami : un médicament soigne, une drogue détruit.
\end{attention}

\begin{propriete}[Les suffixes -los et -igkeit / -heit]
Le suffixe \cle{-los} signifie « sans … » : \all{arbeitslos} (sans travail), \all{obdachlos} (sans logement), \all{hoffnungslos} (sans espoir). Les suffixes \cle{-igkeit} et \cle{-heit} forment des noms abstraits, toujours féminins : \all{arbeitslos} $\rightarrow$ \all{die Arbeitslosigkeit}, \all{krank} $\rightarrow$ \all{die Krankheit}, \all{frei} $\rightarrow$ \all{die Freiheit}. Grâce à ces suffixes, tu peux deviner le sens et le genre de beaucoup de mots nouveaux.
\end{propriete}

\begin{exemplebox}[Les problèmes dans des phrases]
\all{Viele Familien in der Stadt leben in Armut.} « Beaucoup de familles en ville vivent dans la pauvreté. »

\all{Die Arbeitslosigkeit ist ein großes Problem.} « Le chômage est un grand problème. »

\all{Der Hunger ist in manchen Ländern sehr stark.} « La faim est très forte dans certains pays. »

\all{Die Gewalt in den Schulen nimmt zu.} « La violence dans les écoles augmente. » (verbe à particule : \all{zunehmen}, « augmenter » ; la particule \all{zu} va à la fin.)
\end{exemplebox}

\courssub{Champ 2 : les conséquences}

\begin{center}
\begin{tabularx}{\linewidth}{|l|X|X|}
\hline
\rowcolor{popblueL} Deutsch & Pluriel & Français \\
\hline
die Folge & die Folgen & la conséquence \\
\hline
die Krankheit & die Krankheiten & la maladie \\
\hline
die Abhängigkeit & die Abhängigkeiten & la dépendance \\
\hline
der Schulabbruch & die Schulabbrüche & le décrochage scolaire \\
\hline
die Ausgrenzung & die Ausgrenzungen & l'exclusion \\
\hline
\end{tabularx}
\end{center}

\begin{exemplebox}[Les conséquences dans des phrases]
\all{Drogen haben schwere Folgen für die Gesundheit.} « La drogue a de graves conséquences pour la santé. »

\all{Die Abhängigkeit von Drogen ist schwer zu bekämpfen.} « La dépendance à la drogue est difficile à combattre. »

\all{Armut führt oft zum Schulabbruch.} « La pauvreté mène souvent au décrochage scolaire. » (\all{führen zu} + datif : « mener à » ; \all{zum} = \all{zu dem}.)
\end{exemplebox}

\courssub{Champ 3 : les personnes}

\begin{attention}[Les noms de personnes en -e se déclinent]
\all{Der Obdachlose}, \all{der Arbeitslose}, \all{der Süchtige}, \all{der Jugendliche} sont des adjectifs nominalisés : ils se déclinent comme des adjectifs. Après l'article défini : \all{der Arbeitslose}, \all{die Arbeitslosen} (pluriel). Au niveau Première, retiens simplement la forme avec \all{der} et le pluriel en \all{-n}.
\end{attention}

\begin{center}
\begin{tabularx}{\linewidth}{|l|X|X|}
\hline
\rowcolor{popblueL} Deutsch & Pluriel & Français \\
\hline
der Obdachlose & die Obdachlosen & le sans-abri \\
\hline
der Arbeitslose & die Arbeitslosen & le chômeur \\
\hline
der Süchtige & die Süchtigen & le toxicomane, le dépendant \\
\hline
das Opfer & die Opfer & la victime \\
\hline
der Jugendliche & die Jugendlichen & le jeune, l'adolescent \\
\hline
\end{tabularx}
\end{center}

\begin{exemplebox}[Les personnes dans des phrases]
\all{Der Obdachlose schläft auf der Straße.} « Le sans-abri dort dans la rue. »

\all{Viele Jugendliche sind Opfer von Gewalt.} « Beaucoup de jeunes sont victimes de violence. »

\all{Der Süchtige braucht Hilfe.} « Le toxicomane a besoin d'aide. »
\end{exemplebox}

\courssub{Champ 4 : verbes et adjectifs}

\begin{center}
\begin{tabularx}{\linewidth}{|X|X|X|}
\hline
\rowcolor{popblueL} Deutsch & Français & Exemple \\
\hline
nehmen (Drogen nehmen) & prendre (de la drogue) & \all{Er nimmt Drogen.} \\
\hline
bekämpfen & lutter contre & \all{Wir bekämpfen die Armut.} \\
\hline
helfen (+ datif) & aider & \all{Sie hilft dem Obdachlosen.} \\
\hline
leiden an (+ datif) & souffrir de & \all{Er leidet an einer Krankheit.} \\
\hline
gefährlich & dangereux & \all{Drogen sind gefährlich.} \\
\hline
arm & pauvre & \all{Die Familie ist arm.} \\
\hline
süchtig & dépendant & \all{Er ist süchtig nach Drogen.} \\
\hline
schwer & grave, lourd & \all{Das ist ein schweres Problem.} \\
\hline
\end{tabularx}
\end{center}

\begin{attention}[helfen et leiden an : attention au cas]
\all{Helfen} se construit avec le datif : \all{Ich helfe dem Mann.} « J'aide l'homme » (et non \all{den Mann}). \all{Leiden an} se construit aussi avec le datif : \all{Sie leidet an einer Krankheit.} « Elle souffre d'une maladie. » Le verbe \all{nehmen} est fort : \all{er nimmt}, parfait \all{er hat genommen}. Note aussi : \all{süchtig nach} + datif « dépendant de » : \all{Er ist süchtig nach Drogen.}
\end{attention}

\begin{exemplebox}[Expressions utiles]
\all{in Armut leben} « vivre dans la pauvreté » : \all{Viele Menschen leben in Armut.}

\all{gegen Drogen kämpfen} « lutter contre la drogue » : \all{Die Polizei kämpft gegen Drogen.}

\all{die Folgen tragen} « en porter les conséquences » : \all{Die Kinder tragen die Folgen.}
\end{exemplebox}

\courssub{Les familles de mots}

Les familles de mots permettent d'apprendre vite : à partir d'un mot connu, tu construis les autres.

\begin{center}
\begin{tabularx}{\linewidth}{|X|X|X|}
\hline
\rowcolor{popblueL} Base & Dérivé & Français \\
\hline
arm (pauvre) & die Armut & la pauvreté \\
\hline
arbeiten (travailler) & die Arbeit & le travail \\
\hline
die Arbeit & die Arbeitslosigkeit & le chômage \\
\hline
gefährlich (dangereux) & die Gefahr & le danger \\
\hline
folgen (suivre) & die Folge & la conséquence \\
\hline
süchtig (dépendant) & die Sucht & la dépendance, l'addiction \\
\hline
\end{tabularx}
\end{center}

\begin{popfigure}
\begin{tikzpicture}[grow=down, level distance=1.4cm,
  every node/.style={draw=popblue, fill=popblueL, rounded corners, align=center, font=\small, text=popdark},
  edge from parent/.style={draw=poporange, thick, -{Stealth}}]
\node[fill=poporangeL, draw=poporange, font=\bfseries, align=center]{arbeiten\\ (travailler)}
  child { node[align=center]{die Arbeit\\ (le travail)} }
  child { node[align=center]{der Arbeiter\\ (l'ouvrier)} }
  child { node[align=center]{arbeitslos\\ (au chômage)}
    child { node[fill=poppinkL, draw=poppink, align=center]{die Arbeitslosigkeit\\ (le chômage)} }
  };
\end{tikzpicture}
\legende{Arbre de la famille de mots de \all{arbeiten} : une racine, quatre dérivés.}
\end{popfigure}

\begin{savaistu}[Arm, Armut et die Not]
L'adjectif allemand \all{arm} a donné le nom \all{die Armut} grâce au vieux suffixe \all{-ut}. En Suisse alémanique, on emploie aussi souvent \all{die Not} pour parler de la misère : \all{in Not leben} « vivre dans le besoin ». Et savais-tu que le Cameroun et l'Allemagne coopèrent depuis longtemps dans la lutte contre la pauvreté, par exemple dans les projets de formation professionnelle pour les jeunes ?
\end{savaistu}

\courssub{Texte d'application : Ein Tag im Zentrum}

\begin{textebox}[Ein Tag im Zentrum — Texte original]
In Yaoundé gibt es ein kleines Zentrum für Jugendliche. Viele junge Menschen kommen jeden Tag dorthin, weil sie Probleme haben. Manche leben in Armut, andere haben die Schule abgebrochen, und einige sind süchtig nach Drogen.

Frau Mbarga arbeitet im Zentrum. Sie hilft den Jugendlichen. „Die Folgen der Drogen sind sehr schwer“, sagt sie. „Viele Jugendliche leiden an Krankheiten, und ihre Familien tragen auch die Folgen.“

Am Morgen gibt es Essen für die Obdachlosen. Ein junger Mann, Peter, erzählt: „Ich war zwei Jahre arbeitslos. Ich hatte kein Geld und keine Wohnung. Das Zentrum hat mir geholfen. Jetzt lerne ich einen Beruf.“

Die Regierung und viele Organisationen bekämpfen die Armut und die Kriminalität in Kamerun. Aber der Kampf gegen die Drogen ist schwer. „Wir müssen den Jugendlichen eine Zukunft geben“, sagt Frau Mbarga. „Ohne Arbeit und ohne Schule ist die Gefahr groß.“

Am Abend geht Peter nach Hause. Er ist heute zufrieden, denn er hat Hoffnung. Das Zentrum zeigt: Mit Hilfe können Jugendliche die Folgen der sozialen Probleme überwinden.
\end{textebox}

\textbf{Glossaire :} \all{das Zentrum, die Zentren} (le centre) ; \all{abbrechen, bricht ab, hat abgebrochen} (abandonner, interrompre) ; \all{die Wohnung, –en} (le logement) ; \all{der Beruf, –e} (le métier) ; \all{der Kampf, die Kämpfe} (la lutte) ; \all{die Zukunft (–)} (l'avenir) ; \all{die Hoffnung, –en} (l'espoir) ; \all{überwinden, überwindet, hat überwunden} (surmonter).

\textbf{Questions de compréhension :}
\begin{enumerate}
\item Wer arbeitet im Zentrum? (Qui travaille au centre ?)
\item Welche Probleme haben die Jugendlichen? (Quels problèmes ont les jeunes ?)
\item Was hat Peter gelernt? (Qu'est-ce que Peter a appris ?)
\end{enumerate}

\begin{corrige}[Questions sur le texte]
1. \all{Frau Mbarga arbeitet im Zentrum.} 2. \all{Sie leben in Armut, sie haben die Schule abgebrochen, und einige sind süchtig nach Drogen.} 3. \all{Er lernt einen Beruf.}
\end{corrige}

\begin{exoresolu}[Complète avec le bon mot]
Énoncé : complète avec \all{die Armut}, \all{die Arbeitslosigkeit}, \all{die Folgen}, \all{gefährlich} ou \all{der Obdachlose}.

1. \all{Viele Menschen leben in \dots} \quad 2. \all{Drogen sind sehr \dots} \quad 3. \all{\dots\ schläft auf der Straße.} \quad 4. \all{Die Jugendlichen tragen \dots\ der Gewalt.} \quad 5. \all{\dots\ ist ein großes Problem in vielen Ländern.}
\tcblower
Solution : 1. \all{Viele Menschen leben in Armut.} (la pauvreté, nom abstrait féminin, sans article après \all{in}) ; 2. \all{Drogen sind sehr gefährlich.} (adjectif, pas de terminaison après \all{sein}) ; 3. \all{Der Obdachlose schläft auf der Straße.} (sujet masculin singulier, le verbe \all{schläft} le confirme) ; 4. \all{Die Jugendlichen tragen die Folgen der Gewalt.} (pluriel : \all{die Folge} $\rightarrow$ \all{die Folgen}) ; 5. \all{Die Arbeitslosigkeit ist ein großes Problem in vielen Ländern.} (le chômage, féminin en \all{-igkeit}).
\end{exoresolu}

\begin{attention}[Erreurs fréquentes des francophones]
\begin{itemize}
\item N'oublie pas la majuscule : \all{die Armut}, \all{der Hunger}, jamais \all{armut}.
\item Ne traduis pas « les conséquences » par \all{die Konsequenzen} dans tous les contextes : le mot courant est \all{die Folgen}.
\item \all{die Gewalt} est féminin et sans pluriel ; ne dis pas \all{die Gewalte}.
\item « Aider quelqu'un » : \all{jemandem helfen} (datif), pas d'accusatif.
\end{itemize}
\end{attention}

\begin{exercice}[À toi de jouer]
1. Donne l'article et le pluriel (s'il existe) de : \all{Armut}, \all{Folge}, \all{Krankheit}, \all{Hunger}, \all{Opfer}.

2. Forme le nom abstrait : \all{arbeitslos} $\rightarrow$ \dots ; \all{krank} $\rightarrow$ \dots ; \all{süchtig} $\rightarrow$ \dots

3. Traduis en allemand : a) Beaucoup de jeunes souffrent de la pauvreté. b) Le gouvernement lutte contre la drogue. c) Le chômeur a besoin d'aide.
\end{exercice}

\begin{corrige}[Exercice : À toi de jouer]
1. \all{die Armut (–)} ; \all{die Folge (–n)} ; \all{die Krankheit (–en)} ; \all{der Hunger (–)} ; \all{das Opfer (–)}.

2. \all{die Arbeitslosigkeit} ; \all{die Krankheit} ; \all{die Sucht}.

3. a) \all{Viele Jugendliche leiden an Armut.} b) \all{Die Regierung kämpft gegen die Drogen.} (ou : \all{Die Regierung bekämpft die Drogen.}) c) \all{Der Arbeitslose braucht Hilfe.}
\end{corrige}

\begin{aretenir}[L'essentiel du lexique]
\begin{itemize}
\item Les problèmes : \all{die Armut}, \all{die Arbeitslosigkeit}, \all{die Drogen}, \all{die Gewalt}, \all{die Kriminalität}, \all{die Obdachlosigkeit}, \all{der Hunger}.
\item Les conséquences : \all{die Folge (–n)}, \all{die Krankheit (–en)}, \all{die Abhängigkeit (–en)}, \all{der Schulabbruch (–¨e)}, \all{die Ausgrenzung (–en)}.
\item Les personnes : \all{der Obdachlose}, \all{der Arbeitslose}, \all{der Süchtige}, \all{das Opfer}, \all{der Jugendliche}.
\item Verbes-clés : \all{nehmen}, \all{bekämpfen}, \all{helfen (+ Dat.)}, \all{leiden an (+ Dat.)}.
\item Les suffixes \all{-los} et \all{-igkeit/-heit} t'aident à deviner sens et genre ; les noms abstraits comme \all{die Armut} et \all{die Gewalt} n'ont pas de pluriel.
\end{itemize}
\end{aretenir}

\coursec{Grammaire 1 : la cause avec \all{weil} et \all{denn}}

Dans la section précédente, nous avons découvert le vocabulaire des fléaux sociaux : \cle{die Armut} (la pauvreté), \cle{die Drogen} (la drogue), \cle{die Gewalt} (la violence), \cle{die Arbeitslosigkeit} (le chômage). Mais un bon lecteur ne se contente pas de nommer les problèmes : il cherche aussi à les \emph{expliquer}. Pourquoi les gens vivent-ils dans la pauvreté ? Pourquoi certains jeunes prennent-ils de la drogue ? Pour répondre à ces questions en allemand, il faut savoir exprimer la \cle{cause}. C'est exactement ce que font les deux petits mots que nous allons étudier maintenant : \all{weil} et \all{denn}, qui correspondent tous les deux au français « parce que / car ». Attention : ils ne fonctionnent pas du tout de la même manière !

\courssub{Observation : deux phrases du texte support}

Reprenons deux phrases tirées de notre texte \all{Soziale Probleme in unserer Welt} et observons-les attentivement.

\begin{exemplebox}[Deux phrases à observer]
\all{Viele Menschen leben in Armut, weil sie keine Arbeit haben.}\\
« Beaucoup de gens vivent dans la pauvreté parce qu'ils n'ont pas de travail. »

\all{Manche Jugendliche nehmen Drogen, denn sie haben keine Perspektive.}\\
« Certains jeunes prennent de la drogue car ils n'ont pas de perspective. »
\end{exemplebox}

Regardez bien la place du verbe conjugué dans chaque phrase :

\begin{itemize}
\item Dans la première phrase, après \all{weil}, le verbe \all{haben} se trouve \textbf{tout à la fin} : \all{weil sie keine Arbeit HABEN}.
\item Dans la deuxième phrase, après \all{denn}, le verbe \all{haben} reste \textbf{en deuxième position}, comme dans une phrase normale : \all{denn SIE HABEN keine Perspektive}.
\end{itemize}

Ce n'est pas un hasard : c'est la différence fondamentale entre les deux mots. \all{weil} est une \cle{conjonction de subordination} : elle introduit une \emph{subordonnée de cause}, et en allemand, dans toute subordonnée, le verbe conjugué va à la fin. \all{denn} est une \cle{conjonction de coordination} : elle relie deux phrases principales, et la structure ne change pas.

\courssub{La règle de \all{weil} : le verbe à la fin}

\begin{propriete}[La cause avec \all{weil}]
\all{weil} (« parce que ») introduit une \textbf{proposition subordonnée de cause}. Dans cette subordonnée, le \textbf{verbe conjugué se place à la fin} de la proposition.

\all{Er bleibt zu Hause, weil er krank IST.}\\
« Il reste à la maison parce qu'il est malade. »

Structure : [phrase principale] + \all{weil} + [sujet + compléments + \textbf{VERBE à la fin}].
\end{propriete}

Cette règle est absolue et sans exception : dès que \all{weil} apparaît, le verbe conjugué file à la dernière place de la subordonnée. S'il y a un auxiliaire et un participe (au parfait, par exemple), c'est l'\textbf{auxiliaire conjugué} qui va à la fin, juste après le participe :

\begin{exemplebox}[\all{weil} avec différents temps]
\all{Sie nimmt Drogen, weil sie Probleme hat.}\\
« Elle prend de la drogue parce qu'elle a des problèmes. »

\all{Er ist aggressiv, weil er Gewalt zu Hause erlebt hat.}\\
« Il est agressif parce qu'il a vécu la violence à la maison. » (parfait : \all{hat} à la fin, après \all{erlebt})

\all{Die Familie ist arm, weil der Vater seine Arbeit verloren hat.}\\
« La famille est pauvre parce que le père a perdu son travail. »

\all{Viele Kinder gehen nicht zur Schule, weil die Eltern kein Geld haben.}\\
« Beaucoup d'enfants ne vont pas à l'école parce que les parents n'ont pas d'argent. »

\all{Viele Familien in Yaoundé sind arm, weil der Kakao-Preis niedrig ist.}\\
« Beaucoup de familles à Yaoundé sont pauvres parce que le prix du cacao est bas. »
\end{exemplebox}

\begin{attention}[L'erreur numéro 1 des francophones]
En français, « parce que » ne change rien à la phrase : « il est malade » reste « parce qu'il \emph{est} malade ». Le francophone a donc tendance à écrire en allemand : *\all{weil er ist krank}. C'est \textbf{faux} !

Retenez le slogan : \textbf{\all{weil} = verbe à la fin !}

Correct : \all{weil er krank IST.} / Faux : *\all{weil er ist krank.}
\end{attention}

\courssub{La règle de \all{denn} : rien ne bouge}

\begin{propriete}[La cause avec \all{denn}]
\all{denn} (« car ») est une \textbf{conjonction de coordination}. Elle relie deux phrases principales : la structure reste celle d'une phrase principale normale, avec le \textbf{verbe conjugué en position 2}.

\all{Er bleibt zu Hause, denn er IST krank.}\\
« Il reste à la maison car il est malade. »

Structure : [phrase principale] + \all{denn} + [sujet + \textbf{VERBE en position 2} + compléments].
\end{propriete}

Avec \all{denn}, tout est simple : la deuxième partie se construit exactement comme une phrase indépendante. On met une virgule avant \all{denn}.

\begin{exemplebox}[\all{denn} en contexte social]
\all{Er ist arm, denn er hat keine Arbeit.}\\
« Il est pauvre car il n'a pas de travail. »

\all{Die Jugendlichen sind frustriert, denn sie finden keine Jobs.}\\
« Les jeunes sont frustrés car ils ne trouvent pas d'emploi. »

\all{Sie kann die Schulgebühren nicht bezahlen, denn ihre Familie ist sehr arm.}\\
« Elle ne peut pas payer les frais scolaires car sa famille est très pauvre. »

\all{Der Bauer hat keine Ernte, denn die Dürre war zu stark.}\\
« Le paysan n'a pas de récolte car la sécheresse était trop forte. »
\end{exemplebox}

\courssub{La place dans la phrase : \all{weil} peut commencer, \all{denn} jamais}

Voici la deuxième grande différence. La subordonnée introduite par \all{weil} peut se placer \textbf{en tête de phrase}. Dans ce cas, la subordonnée entière occupe la position 1, et le verbe de la phrase principale vient \textbf{immédiatement après la virgule} (position 2) :

\all{Weil er keine Arbeit hat, IST er arm.}\\
« Parce qu'il n'a pas de travail, il est pauvre. »

\all{Weil sie arm sind, KAUFEN sie keine Bücher.}\\
« Parce qu'ils sont pauvres, ils n'achètent pas de livres. »

En revanche, \all{denn} ne peut \textbf{jamais} commencer une phrase. On ne peut pas dire *\all{Denn er hat keine Arbeit, ist er arm.}

\begin{attention}[\all{denn} jamais en tête de phrase]
En tête de phrase, seul \all{weil} est possible. Et attention au piège : après une subordonnée en \all{weil} placée au début, le verbe de la principale arrive \textbf{juste après la virgule}, avant le sujet :

\all{Weil er krank ist, BLEIBT er zu Hause.} (et non *\all{Weil er krank ist, er bleibt...})

La subordonnée compte comme la position 1 ; le verbe de la principale doit donc être en position 2.
\end{attention}

\begin{aretenir}[Répondre à la question \all{Warum...?}]
À la question \all{Warum?} (« Pourquoi ? »), on peut répondre avec \all{weil} ou avec \all{denn} :

\all{Warum ist er arm? — Weil er keine Arbeit hat.} ou \all{Er ist arm, denn er hat keine Arbeit.}\\
« Pourquoi est-il pauvre ? — Parce qu'il n'a pas de travail. / Il est pauvre car il n'a pas de travail. »

Dans une réponse courte commençant directement par \all{weil}, le verbe va bien sûr à la fin : \all{Weil er keine Arbeit HAT.}
\end{aretenir}

\courssub{Schéma : la place du verbe avec \all{weil} et \all{denn}}

\begin{popfigure}
\begin{tikzpicture}[font=\small]
% weil
\node[draw=popblue, fill=popblueL, rounded corners, minimum width=3.4cm, minimum height=0.9cm, align=center] (h1) at (0,2.2) {Hauptsatz\\ \all{Er bleibt zu Hause,}};
\node[draw=poporange, fill=poporangeL, rounded corners, minimum width=1.1cm, minimum height=0.9cm] (w) at (3.1,2.2) {\all{weil}};
\node[draw=popgreen, fill=popgreenL, rounded corners, minimum width=2.6cm, minimum height=0.9cm, align=center] (s1) at (5.6,2.2) {Sujet + Reste\\ \all{er krank}};
\node[draw=poppink, fill=poppinkL, rounded corners, minimum width=1.7cm, minimum height=0.9cm, align=center] (v1) at (8.2,2.2) {\textbf{VERBE}\\ \all{IST.}};
\draw[-{Stealth}, thick, popdark] (h1) -- (w);
\draw[-{Stealth}, thick, popdark] (w) -- (s1);
\draw[-{Stealth}, thick, popdark] (s1) -- (v1);
% denn
\node[draw=popblue, fill=popblueL, rounded corners, minimum width=3.4cm, minimum height=0.9cm, align=center] (h2) at (0,0) {Hauptsatz\\ \all{Er bleibt zu Hause,}};
\node[draw=poporange, fill=poporangeL, rounded corners, minimum width=1.1cm, minimum height=0.9cm] (d) at (3.1,0) {\all{denn}};
\node[draw=poppink, fill=poppinkL, rounded corners, minimum width=2.6cm, minimum height=0.9cm, align=center] (v2) at (5.6,0) {Sujet + \textbf{VERBE}\\ \all{er IST}};
\node[draw=popgreen, fill=popgreenL, rounded corners, minimum width=1.7cm, minimum height=0.9cm, align=center] (r2) at (8.2,0) {Reste\\ \all{krank.}};
\draw[-{Stealth}, thick, popdark] (h2) -- (d);
\draw[-{Stealth}, thick, popdark] (d) -- (v2);
\draw[-{Stealth}, thick, popdark] (v2) -- (r2);
\end{tikzpicture}
\legende{Avec \all{weil}, le verbe conjugué va à la fin de la subordonnée ; avec \all{denn}, il reste en position 2, comme dans une phrase principale.}
\end{popfigure}

\courssub{Tableau récapitulatif : \all{weil} ou \all{denn} ?}

\begin{center}
\begin{tabularx}{\linewidth}{|X|X|X|}
\hline
\rowcolor{popblueL} & \all{weil} & \all{denn} \\
\hline
Type de mot & conjonction de subordination & conjonction de coordination \\
\hline
Type de proposition & subordonnée de cause & deuxième phrase principale \\
\hline
Place du verbe conjugué & à la \textbf{fin} : \all{weil er krank ist} & en \textbf{position 2} : \all{denn er ist krank} \\
\hline
Peut commencer la phrase ? & \textbf{oui} : \all{Weil er krank ist, bleibt er zu Hause.} & \textbf{jamais} \\
\hline
Équivalent français & parce que & car \\
\hline
Registre & courant, oral et écrit & plutôt écrit, un peu soutenu \\
\hline
\end{tabularx}
\end{center}

\begin{center}
\begin{tabularx}{\linewidth}{|X|X|X|}
\hline
\rowcolor{popblueL} Français & Deutsch & Remarque \\
\hline
Il reste à la maison parce qu'il est malade. & \all{Er bleibt zu Hause, weil er krank ist.} & verbe à la fin après \all{weil} \\
\hline
Il reste à la maison car il est malade. & \all{Er bleibt zu Hause, denn er ist krank.} & structure inchangée après \all{denn} \\
\hline
Parce qu'il est malade, il reste à la maison. & \all{Weil er krank ist, bleibt er zu Hause.} & verbe de la principale juste après la virgule \\
\hline
\end{tabularx}
\end{center}

Remarquez bien : en français, la structure de la phrase ne change \emph{jamais}, que l'on utilise « parce que » ou « car ». En allemand, le choix entre \all{weil} et \all{denn} a des conséquences directes sur la place du verbe. C'est toute la difficulté — et toute la beauté — de la syntaxe allemande !

\courssub{Exercice résolu et entraînement}

\begin{exoresolu}[Transformation : de \all{denn} à \all{weil}]
Énoncé : transformez la phrase suivante en utilisant \all{weil} à la place de \all{denn}.

\all{Manche Jugendliche nehmen Drogen, denn sie haben keine Perspektive.}
\tcblower
Solution détaillée :

Étape 1 : on remplace \all{denn} par \all{weil}. Étape 2 : \all{weil} exige le verbe conjugué à la fin, donc \all{haben} quitte la position 2 et va après \all{keine Perspektive}.

\all{Manche Jugendliche nehmen Drogen, weil sie keine Perspektive haben.}\\
« Certains jeunes prennent de la drogue parce qu'ils n'ont pas de perspective. »

Variante avec \all{weil} en tête
\end{exoresolu}


\coursec{Grammaire 2 : le comparatif — Conjugaison : le parfait}

Dans la section précédente, nous avons appris à exprimer la \textbf{cause} avec \all{weil} et \all{denn} : \all{Viele Jugendliche nehmen Drogen, weil sie Probleme haben.} (Beaucoup de jeunes prennent de la drogue parce qu'ils ont des problèmes.) Nous allons maintenant apprendre à \textbf{comparer} les fléaux sociaux entre eux (\emph{la pauvreté est-elle pire que la faim ?}) grâce au \cle{comparatif}, puis à \textbf{raconter au passé} ce qui s'est passé grâce au \cle{parfait} (\emph{Perfekt}), le temps le plus utilisé de la conversation allemande.

\courssub{Le comparatif : observer avant de comprendre}

\begin{textebox}[Mini-texte d'observation : Vergleiche]
\all{Die Arbeitslosigkeit ist heute größer als früher. Drogen sind gefährlicher als Alkohol. Arme Familien sind oft ärmer als früher, und reiche Familien sind reicher. Viele sagen: Gewalt ist schlimmer als Armut. Aber ist das wirklich so? Für viele junge Menschen in Yaoundé ist die Arbeitslosigkeit das größte Problem. Sie ist wichtiger als alles andere.}
\par\medskip
\emph{Traduction : Le chômage est aujourd'hui plus grand qu'autrefois. La drogue est plus dangereuse que l'alcool. Les familles pauvres sont souvent plus pauvres qu'avant, et les familles riches sont plus riches. Beaucoup disent : la violence est pire que la pauvreté. Mais est-ce vraiment le cas ? Pour beaucoup de jeunes à Yaoundé, le chômage est le plus grand problème. Il est plus important que tout le reste.}
\par\medskip
\textbf{Glossaire :} \all{früher} = autrefois ; \all{reich} = riche ; \all{schlimm} = grave, mauvais ; \all{wichtig} = important ; \all{als} = que (dans une comparaison).
\end{textebox}

\textbf{Observons.} Soulignons les mots qui expriment la comparaison :

\begin{itemize}
\item \all{größ\textbf{er} als früher} — plus grand \textbf{qu'}autrefois ;
\item \all{gefährlich\textbf{er} als Alkohol} — plus dangereux \textbf{que} l'alcool ;
\item \all{ärm\textbf{er} als früher} — plus pauvre \textbf{qu'}avant ;
\item \all{schlimm\textbf{er} als Armut} — pire \textbf{que} la pauvreté.
\end{itemize}

On remarque deux choses : l'adjectif prend la terminaison \all{-er}, et le second terme de la comparaison est introduit par \all{als}. C'est exactement le mécanisme du comparatif allemand.

\courssub{Formation du comparatif}

\begin{propriete}[Règle du comparatif]
Le comparatif se forme en ajoutant \cle{-er} à l'adjectif. Le second terme de la comparaison est introduit par \cle{als} (= « que »).
\par
\all{Die Stadt ist größ\textbf{er} als das Dorf.} (La ville est plus grande que le village.)
\par
\all{Drogen sind gefährlich\textbf{er} als Alkohol.} (La drogue est plus dangereuse que l'alcool.)
\end{propriete}

\textbf{Premier cas particulier : l'Umlaut.} Les adjectifs \textbf{monosyllabiques} (une seule syllabe) dont la voyelle est \all{a}, \all{o} ou \all{u} prennent en général un \textbf{Umlaut} au comparatif :

\begin{center}
\begin{tabular}{|l|l|l|}
\hline
\rowcolor{popblueL} Positif & Comparatif & Traduction \\ \hline
arm & ärmer & pauvre / plus pauvre \\ \hline
groß & größer & grand / plus grand \\ \hline
jung & jünger & jeune / plus jeune \\ \hline
alt & älter & vieux / plus vieux \\ \hline
stark & stärker & fort / plus fort \\ \hline
lang & länger & long / plus long \\ \hline
\end{tabular}
\end{center}

\begin{attention}[Pas d'Umlaut partout !]
Les adjectifs de \textbf{plusieurs syllabes} ne prennent jamais d'Umlaut : \all{gefährlich} $\rightarrow$ \all{gefährlicher}, \all{wichtig} $\rightarrow$ \all{wichtiger}, \all{interessant} $\rightarrow$ \all{interessanter}. De même, quelques monosyllabes échappent à l'Umlaut : \all{falsch} $\rightarrow$ \all{falscher}, \all{klar} $\rightarrow$ \all{klarer}, \all{voll} $\rightarrow$ \all{voller}, \all{stolz} $\rightarrow$ \all{stolzer}. En cas de doute, apprends l'adjectif avec son comparatif, comme un verbe fort.
\end{attention}

\textbf{Deuxième cas particulier : les comparatifs irréguliers.} Quatre formes sont à mémoriser absolument, car elles reviennent sans cesse :

\begin{center}
\begin{tabular}{|l|l|l|}
\hline
\rowcolor{popblueL} Positif & Comparatif & Traduction \\ \hline
gut & besser & bon / meilleur \\ \hline
viel & mehr & beaucoup / plus \\ \hline
gern & lieber & volontiers / plutôt, de préférence \\ \hline
hoch & höher & haut / plus haut \\ \hline
\end{tabular}
\end{center}

Note bien que \all{hoch} perd son \all{-c-} : \all{hoch} $\rightarrow$ \all{höher}. Et comme en français (« bon $\rightarrow$ meilleur », et non *« plus bon »), l'allemand dit \all{gut} $\rightarrow$ \all{besser}, jamais *\all{güter}.

\begin{popfigure}
\begin{center}
\begin{tikzpicture}
\node[draw=popblue, fill=popblueL, rounded corners, font=\bfseries, minimum width=2.6cm] (titre) at (0,0) {Comparatif};
\node[draw=popgreen, fill=popgreenL, rounded corners, align=center] (reg) at (-4.2,-1.8) {Règle générale\\ adjectif + -er\\ \all{gefährlich $\rightarrow$ gefährlicher}};
\node[draw=poporange, fill=poporangeL, rounded corners, align=center] (uml) at (0,-1.8) {Umlaut (a, o, u)\\ monosyllabes\\ \all{arm $\rightarrow$ ärmer}};
\node[draw=poppurple, fill=poppurpleL, rounded corners, align=center] (irr) at (4.2,-1.8) {Irréguliers\\ \all{gut $\rightarrow$ besser}\\ \all{viel $\rightarrow$ mehr}};
\draw[-{Stealth}, thick, popdark] (titre) -- (reg);
\draw[-{Stealth}, thick, popdark] (titre) -- (uml);
\draw[-{Stealth}, thick, popdark] (titre) -- (irr);
\end{tikzpicture}
\end{center}
\legende{Carte mentale : les trois cas de formation du comparatif.}
\end{popfigure}

\courssub{La structure de comparaison : als ou wie ?}

\begin{propriete}[als = inégalité, wie = égalité]
\begin{itemize}
\item \textbf{Inégalité} (« plus... que ») : comparatif + \cle{als}. \all{Er ist älter als ich.} (Il est plus âgé que moi.)
\item \textbf{Égalité} (« aussi... que ») : \cle{so} + adjectif au positif + \cle{wie}. \all{Er ist so alt wie ich.} (Il est aussi âgé que moi.)
\end{itemize}
\end{propriete}

\begin{attention}[Le piège n° 1 des francophones : *größer wie]
En français, on dit « plus grand \emph{que} » et « aussi grand \emph{que} » : le même mot « que » sert aux deux. En allemand, il faut choisir : \all{als} après un comparatif, \all{wie} après \all{so}. La phrase *\all{Die Stadt ist größer wie das Dorf} est une faute classique, immédiatement repérée par le correcteur. Retiens le couple : \textbf{comparatif + als} / \textbf{so + wie}.
\end{attention}

\begin{exemplebox}[Exemples traduits]
\begin{itemize}
\item \all{Armut ist schlimmer als Hunger.} (La pauvreté est pire que la faim.)
\item \all{Die Arbeitslosigkeit ist heute größer als früher.} (Le chômage est aujourd'hui plus important qu'autrefois.)
\item \all{Mein Bruder ist jünger als ich.} (Mon frère est plus jeune que moi.)
\item \all{Fußball ist so populär wie Handball.} (Le football est aussi populaire que le handball.)
\item \all{Ich trinke lieber Wasser als Alkohol.} (Je bois de préférence de l'eau plutôt que de l'alcool.)
\end{itemize}
\end{exemplebox}

\courssub{Comparatif attribut et comparatif épithète}

Jusqu'ici, le comparatif était \textbf{attribut}, placé après \all{sein} ou \all{werden} : \all{Die Gewalt ist gefährlicher.} Dans ce cas, il reste \textbf{invariable}, comme en français (« la violence est plus dangereuse »).

Mais quand le comparatif est \textbf{épithète}, c'est-à-dire placé \textbf{devant le nom}, il prend en plus la \textbf{désinence adjectivale} normale (après \all{ein}, \all{der}, etc.) : la terminaison \all{-er} du comparatif \emph{ne remplace pas} la désinence, elle s'y ajoute.

\begin{center}
\begin{tabular}{|l|l|}
\hline
\rowcolor{popblueL} Attribut (invariable) & Épithète (décliné) \\ \hline
\all{Das Problem ist gefährlicher.} & \all{ein gefährlicher\textbf{es} Problem} \\ \hline
\all{Die Stadt ist größer.} & \all{eine größer\textbf{e} Stadt} \\ \hline
\all{Der Mann ist älter.} & \all{ein älter\textbf{er} Mann} \\ \hline
\end{tabular}
\end{center}

\begin{exemplebox}[Exemples traduits]
\begin{itemize}
\item \all{Drogen sind ein gefährlicheres Problem als Alkohol.} (La drogue est un problème plus dangereux que l'alcool.)
\item \all{Wir brauchen eine bessere Lösung.} (Nous avons besoin d'une meilleure solution.)
\item \all{Ein älterer Mann hilft den Jugendlichen.} (Un homme plus âgé aide les jeunes.)
\end{itemize}
\end{exemplebox}

\begin{exoresolu}[Mets les adjectifs au comparatif]
Énoncé : complète avec le comparatif de l'adjectif entre parenthèses.
\begin{enumerate}
\item \all{Die Gewalt ist \dots\ als früher. (schlimm)}
\item \all{Mein Dorf ist \dots\ als Douala. (klein)}
\item \all{Ein Fahrrad ist \dots\ als ein Auto. (billig)}
\item \all{Wasser ist \dots\ als Alkohol. (gut)}
\item \all{Die Jugendlichen sind \dots\ als ihre Eltern. (jung)}
\end{enumerate}
\tcblower
Solution détaillée :
\begin{enumerate}
\item \all{schlimmer} : adjectif polysyllabique, simple ajout de \all{-er}. « La violence est pire qu'autrefois. »
\item \all{kleiner} : pas d'Umlaut (la voyelle \all{ei} n'est pas concernée). « Mon village est plus petit que Douala. »
\item \all{billiger} : deux syllabes, donc pas d'Umlaut. « Un vélo est moins cher qu'une voiture. »
\item \all{besser} : comparatif irrégulier de \all{gut}. « L'eau est meilleure que l'alcool. »
\item \all{jünger} : monosyllabe en \all{u} $\rightarrow$ Umlaut. « Les jeunes sont plus jeunes que leurs parents. »
\end{enumerate}
\end{exoresolu}

\courssub{Conjugaison : le parfait (Perfekt)}

\textbf{Observons d'abord} trois phrases tirées de notre thème :

\begin{itemize}
\item \all{Die Gewalt \textbf{hat} zuge\textbf{nommen}.} (La violence a augmenté.)
\item \all{Viele junge Menschen \textbf{haben} die Schule ver\textbf{lassen}.} (Beaucoup de jeunes ont quitté l'école.)
\item \all{Viele Menschen \textbf{sind} in die Stadt ge\textbf{zogen}.} (Beaucoup de gens ont emménagé en ville.)
\end{itemize}

Dans chaque phrase, on trouve \textbf{deux} éléments verbaux : un auxiliaire conjugué (\all{hat}, \all{haben}, \all{sind}) en \textbf{deuxième position}, et un participe passé (\all{zugenommen}, \all{verlassen}, \all{gezogen}) tout \textbf{à la fin}. C'est le \cle{parfait} allemand, l'équivalent du passé composé français.

\begin{propriete}[Formation du parfait]
Parfait = auxiliaire \cle{haben} ou \cle{sein} conjugué au présent (position 2) + \cle{Partizip II} (participe passé) en \textbf{fin de phrase}. C'est le fameux « cadre verbal » allemand : tout le reste de la phrase se loge entre les deux.
\end{propriete}

\begin{popfigure}
\begin{center}
\begin{tikzpicture}
\node[draw=popblue, fill=popblueL, rounded corners, align=center, minimum width=2.4cm] (s) at (0,0) {Sujet\\ \all{Viele Menschen}};
\node[draw=poporange, fill=poporangeL, rounded corners, align=center, minimum width=2.4cm] (aux) at (3.4,0) {Auxiliaire (pos. 2)\\ \all{sind}};
\node[draw=popgreen, fill=popgreenL, rounded corners, align=center, minimum width=2.8cm] (mil) at (7.0,0) {Compléments\\ \all{in die Stadt}};
\node[draw=poppurple, fill=poppurpleL, rounded corners, align=center, minimum width=2.8cm] (p2) at (10.8,0) {Partizip II (fin)\\ \all{gezogen}};
\draw[-{Stealth}, thick, popdark] (s) -- (aux);
\draw[-{Stealth}, thick, popdark] (aux) -- (mil);
\draw[-{Stealth}, thick, popdark] (mil) -- (p2);
\draw[very thick, poppink] (aux.south west) -- (p2.south east) node[midway, below=4pt, poppink, align=center]{cadre verbal};
\end{tikzpicture}
\end{center}
\legende{Le cadre verbal du parfait : l'auxiliaire ouvre, le Partizip II ferme.}
\end{popfigure}

\courssub{Haben ou sein ? Le choix de l'auxiliaire}

\begin{propriete}[Choix de l'auxiliaire]
\begin{itemize}
\item \textbf{haben} : c'est le cas général (la grande majorité des verbes) : \all{Ich habe gearbeitet.} (J'ai travaillé.)
\item \textbf{sein} : pour les verbes de \textbf{déplacement} (\all{gehen, kommen, fahren, ziehen}) et de \textbf{changement d'état} (\all{sterben, aufwachen, werden}), ainsi que \all{bleiben} et \all{sein} : \all{Er ist gekommen.} (Il est venu.) \all{Sie ist gestorben.} (Elle est morte.)
\end{itemize}
\end{propriete}

\begin{attention}[Comparaison avec le français]
Le passé composé français ressemble beaucoup au parfait allemand (« j'ai pris » = \all{ich habe genommen}), mais attention : le français utilise « être » seulement avec une petite liste de verbes (aller, venir, partir...), alors que l'allemand utilise \all{sein} avec \textbf{tous} les verbes de déplacement et de changement d'état. Deux cas typiques où les deux langues coïncident heureusement : \all{Er ist geblieben.} (Il est resté.) et \all{Er ist gestorben.} (Il est mort.) Mais méfie-toi : \all{fliegen} (voler, partir en avion) prend aussi \all{sein} : \all{Sie ist nach Deutschland geflogen.} (Elle est partie en Allemagne en avion.)
\end{attention}

\courssub{Le Partizip II : faible, fort, mixte}

\begin{center}
\begin{tabular}{|l|l|l|l|}
\hline
\rowcolor{popblueL} Type & Formation & Exemple & Traduction \\ \hline
Faible & ge + radical + t & \all{machen $\rightarrow$ gemacht} & fait \\ \hline
Faible (radical en -t, -d) & ge + radical + et & \all{arbeiten $\rightarrow$ gearbeitet} & travaillé \\ \hline
Verbe en -ieren & pas de ge-, fin en -t & \all{studieren $\rightarrow$ studiert} & étudié \\ \hline
Fort & ge + (voyelle changée) + en & \all{nehmen $\rightarrow$ genommen} & pris \\ \hline
Fort & idem & \all{bleiben $\rightarrow$ geblieben} & resté \\ \hline
Préfixe inséparable & pas de ge- & \all{verlassen $\rightarrow$ verlassen} & quitté \\ \hline
Préfixe séparable & ge- au milieu & \all{zunehmen $\rightarrow$ zugenommen} & augmenté \\ \hline
\end{tabular}
\end{center}

Trois réflexes à acquérir :
\begin{enumerate}
\item \textbf{Verbes faibles} : \all{ge-...-t} : \all{machen $\rightarrow$ gemacht}, \all{lernen $\rightarrow$ gelernt}. Si le radical finit en \all{-t} ou \all{-d}, on ajoute un \all{-e-} de liaison : \all{arbeiten $\rightarrow$ gearbeitet}, \all{warten $\rightarrow$ gewartet}. Les verbes en \all{-ieren} (souvent d'origine française) n'ont \textbf{pas} de \all{ge-} : \all{studieren $\rightarrow$ studiert}, \all{diskutieren $\rightarrow$ diskutiert}.
\item \textbf{Verbes forts} : \all{ge-...-en}, souvent avec changement de voyelle : \all{nehmen $\rightarrow$ genommen}, \all{gehen $\rightarrow$ gegangen}, \all{trinken $\rightarrow$ getrunken}. Ils s'apprennent par cœur, comme les verbes irréguliers anglais.
\item \textbf{Préfixes} : inséparables (\all{ver-, be-, er-}) $\rightarrow$ pas de \all{ge-} : \all{verlassen $\rightarrow$ verlassen}, \all{bekommen $\rightarrow$ bekommen} ; séparables $\rightarrow$ \all{ge-} inséré entre la particule et le verbe : \all{zunehmen $\rightarrow$ zuge\textbf{nommen}}, \all{mitmachen $\rightarrow$ mitge\textbf{macht}}.
\end{enumerate}

\begin{exemplebox}[Exemples au parfait, traduits]
\begin{itemize}
\item \all{Er hat gestern Drogen genommen.} (Il a pris de la drogue hier.)
\item \all{Viele Menschen sind in die Stadt gezogen.} (Beaucoup de gens ont emménagé en ville.)
\item \all{Die Gewalt hat zugenommen.} (La violence a augmenté.)
\item \all{Wir haben viel über Armut gelernt.} (Nous avons beaucoup appris sur la pauvreté.)
\item \all{Sie ist nach Deutschland gefahren.} (Elle est partie en Allemagne.)
\end{itemize}
\end{exemplebox}

\begin{attention}[Le Partizip II reste à la fin, quoi qu'il arrive]
Au parfait, le participe se place \textbf{toujours} en fin de phrase, même si la phrase est longue : \all{Er hat letztes Jahr wegen der Arbeitslosigkeit die Schule verlassen.} (Il a quitté l'école l'année dernière à cause du chômage.) Le francophone a tendance à rapprocher le participe de l'auxiliaire, comme en français (« il a quitté l'école ») : résiste à cette tentation ! L'allemand « emballe » toute la phrase dans le cadre verbal.
\end{attention}

\begin{exoresolu}[Mets les phrases au parfait]
Énoncé : récris au parfait.
\begin{enumerate}
\item \all{Die Jugendlichen lernen viel über Drogen.}
\item \all{Er geht in die Stadt.}
\item \all{Viele Familien verlieren ihr Geld.}
\end{enumerate}
\tcblower
Solution détaillée :
\begin{enumerate}
\item \all{Die Jugendlichen haben viel über Drogen gelernt.} — verbe faible, auxiliaire \all{haben}. (Les jeunes ont beaucoup appris sur la drogue.)
\item \all{Er ist in die Stadt gegangen.} — verbe fort de déplacement : auxiliaire \all{sein}, participe irrégulier \all{gegangen}. (Il est allé en ville.)
\item \all{Viele Familien haben ihr Geld verloren.} — \all{verlieren} est fort (\all{verloren}) et son préfixe \all{ver-} est inséparable : pas de \all{ge-}. (Beaucoup de familles ont perdu leur argent.)
\end{enumerate}
\end{exoresolu}

\begin{aretenir}[L'essentiel de la section]
\begin{itemize}
\item Comparatif = adjectif + \all{-er} ; Umlaut pour les monosyllabes en a/o/u (\all{arm $\rightarrow$ ärmer}) ; irréguliers : \all{gut $\rightarrow$ besser}, \all{viel $\rightarrow$ mehr}, \all{gern $\rightarrow$ lieber}, \all{hoch $\rightarrow$ höher}.
\item Inégalité : comparatif + \all{als} ; égalité : \all{so ... wie}. Jamais *\all{größer wie}.
\item Devant un nom, le comparatif se décline : \all{ein gefährlicheres Problem}, \all{eine bessere Lösung}.
\item Parfait = \all{haben}/\all{sein} (position 2) + Partizip II (fin de phrase) : le cadre verbal.
\item \all{sein} avec les verbes de déplacement et de changement d'état (\all{gehen, kommen, ziehen, sterben, bleiben}).
\item Partizip II : faible \all{gemacht}, fort \all{genommen}, verbe en \all{-ieren} \all{studiert}, préfixe inséparable \all{verlassen}, séparable \all{zugenommen}.
\end{itemize}
\end{aretenir}

\begin{savaistu}[Le parfait, roi de la conversation]
Au Bac comme dans la vraie vie, le parfait est le temps le plus fréquent pour raconter au passé à l'oral et dans les textes courants (lettres, dialogues, journaux intimes). Le prétérit (\all{ich ging, ich machte}) est surtout réservé à l'écrit littéraire et aux verbes \all{sein} (\all{war}), \all{haben} (\all{hatte}) et aux verbes modaux (\all{musste, konnte, wollte}). Un Allemand dira naturellement : \all{Ich bin gestern in die Stadt gegangen} et non \all{Ich ging gestern in die Stadt}. Pour tes productions écrites et orales de Première, le parfait est donc ton meilleur allié !
\end{savaistu}

\coursec{Méthode du commentaire et production guidée}

Dans la section précédente, nous avons appris à comparer (\all{größer als}, \all{gefährlicher als}) et à raconter au passé avec le parfait (\all{Viele Familien haben ihr Haus verloren}). Ces outils ne servent pas seulement à traduire : ils servent aussi à \cle{parler d'un texte} et à \cle{écrire un court paragraphe}. C'est exactement ce que l'on te demandera à l'examen : lire un texte sur un thème social, le commenter, puis produire quelques phrases personnelles. Cette dernière section te donne la méthode complète, étape par étape, avec un modèle rédigé et une production guidée.

\courssub{La méthode du commentaire de texte en 4 étapes}

\courssub{Observer d'abord : un mini-commentaire}

Avant d'apprendre la méthode, observons comment un élève de Première peut parler d'un texte en allemand. Lis ces quatre phrases :

\begin{exemplebox}[Observer avant de comprendre]
\all{Der Text handelt von den sozialen Problemen in der Welt.} — Le texte traite des problèmes sociaux dans le monde.\par
\all{Die Hauptidee ist: Armut, Drogen und Gewalt haben gefährliche Folgen.} — L'idée principale est : la pauvreté, la drogue et la violence ont des conséquences dangereuses.\par
\all{Der Autor zeigt, dass viele Jugendliche keine Arbeit finden.} — L'auteur montre que beaucoup de jeunes ne trouvent pas de travail.\par
\all{Meiner Meinung nach ist die Arbeitslosigkeit das größte Problem.} — À mon avis, le chômage est le plus grand problème.
\end{exemplebox}

Observe bien : chaque phrase remplit une \cle{fonction} différente. La première \textbf{présente} le texte, la deuxième donne l'\textbf{idée principale}, la troisième \textbf{analyse} une idée secondaire, la quatrième donne un \textbf{avis personnel}. Un commentaire de texte, même court, suit toujours ce chemin en quatre étapes.

\begin{methode}[Commenter un texte en 4 étapes]
\begin{enumerate}
\item \textbf{Présenter le texte} : dis de quel type de texte il s'agit (un article, un reportage, une interview), quel est son thème et d'où il vient (la source). Phrase-outil : \all{Der Text handelt von...} (Le texte traite de...).
\item \textbf{Dégager l'idée principale} : en une ou deux phrases, dis ce que le texte veut dire essentiellement. Phrase-outil : \all{Die Hauptidee ist...} (L'idée principale est...) ou \all{Der Autor zeigt, dass...} (L'auteur montre que...).
\item \textbf{Analyser les idées secondaires et le vocabulaire} : relève deux ou trois informations précises (chiffres, exemples, causes, conséquences) et commente les mots importants du thème (\all{die Armut}, \all{die Drogen}, \all{die Folge}...).
\item \textbf{Donner ton avis argumenté} : dis ce que tu penses et justifie avec \all{weil} ou \all{denn}. Phrases-outils : \all{Ich finde, dass...} (Je trouve que...) ou \all{Meiner Meinung nach...} (À mon avis...).
\end{enumerate}
\end{methode}

\begin{popfigure}
\begin{tikzpicture}[
  node distance=0.7cm,
  etape/.style={rectangle, rounded corners=4pt, draw=popblue, fill=popblueL, thick, minimum width=6.5cm, minimum height=1.1cm, align=center, font=\small},
  fleche/.style={-{Stealth[length=3mm]}, thick, poporange}
]
\node[etape, align=center] (e1){\textbf{1. Présentation}\\ \all{Der Text handelt von...}};
\node[etape, below=of e1, align=center] (e2){\textbf{2. Idée principale}\\ \all{Der Autor zeigt, dass...}};
\node[etape, below=of e2, align=center] (e3){\textbf{3. Analyse}\\ idées secondaires + vocabulaire};
\node[etape, below=of e3, align=center] (e4){\textbf{4. Avis personnel}\\ \all{Meiner Meinung nach...}};
\draw[fleche] (e1) -- (e2);
\draw[fleche] (e2) -- (e3);
\draw[fleche] (e3) -- (e4);
\end{tikzpicture}
\legende{Organigramme du commentaire de texte : quatre étapes dans l'ordre, de la présentation à l'avis personnel.}
\end{popfigure}

\begin{attention}[Après \all{dass}, le verbe va à la fin !]
Pour donner ton avis, utilise \all{Ich finde, dass...} ou \all{Meiner Meinung nach...}. Mais attention, piège classique du francophone : après \all{dass}, le verbe conjugué se place \textbf{à la fin} de la phrase, comme après \all{weil}.\par
\all{Ich finde, dass die Armut ein großes Problem \textbf{ist}.} — Je trouve que la pauvreté est un grand problème.\par
Faux : \emph{Ich finde, dass die Armut ist ein großes Problem.} (ordre des mots français !)\par
Astuce : \all{dass} = \all{weil} = verbe à la fin. \all{denn} = verbe en deuxième position.
\end{attention}

\courssub{Un commentaire modèle rédigé}

Voici un commentaire complet du texte support \all{Soziale Probleme in unserer Welt}, rédigé par un élève. Les étapes de la méthode sont signalées entre crochets : repère-les et réutilise ce modèle comme gabarit.

\begin{textebox}[Kommentar — commentaire modèle]
[1. Présentation] Der Text handelt von den sozialen Problemen in unserer Welt. Es ist ein kurzer Artikel über Armut, Drogen, Gewalt und Arbeitslosigkeit.\par
[2. Idée principale] Der Autor zeigt, dass diese Probleme gefährliche Folgen für die Menschen und besonders für die Jugendlichen haben.\par
[3. Analyse] Der Text erklärt: Viele Menschen sind arm, weil sie keine Arbeit haben. Die Folge ist, dass viele Jugendliche die Schule verlassen. Das Wort „Folge“ ist sehr wichtig: Jedes Problem hat eine Konsequenz.\par
[4. Avis personnel] Ich finde, dass die Regierung mehr helfen muss, denn die Jugendlichen sind die Zukunft. Meiner Meinung nach ist Bildung die beste Lösung.
\end{textebox}

\textbf{Glossaire :} \all{der Artikel, die Artikel} (l'article) ; \all{die Folge, -n} (la conséquence) ; \all{die Jugendlichen} (les jeunes, pluriel seulement) ; \all{die Regierung, -en} (le gouvernement) ; \all{die Lösung, -en} (la solution) ; \all{helfen, half, hat geholfen} (aider, irrégulier) ; \all{verlassen, verließ, hat verlassen} (quitter, irrégulier).

Remarque comme ce commentaire réutilise \textbf{toute la leçon} : une cause avec \all{weil} (verbe à la fin), une cause avec \all{denn} (verbe en position 2), le vocabulaire du thème et les phrases-outils. C'est exactement ce que le correcteur attend.

\begin{exemplebox}[Les phrases-outils du commentaire]
\all{Der Text handelt von der Arbeitslosigkeit.} — Le texte traite du chômage.\par
\all{Der Autor zeigt, dass Drogen sehr gefährlich sind.} — L'auteur montre que la drogue est très dangereuse.\par
\all{Ich finde, dass Gewalt in den Schulen ein großes Problem ist.} — Je trouve que la violence dans les écoles est un grand problème.\par
\all{Meiner Meinung nach ist die Arbeitslosigkeit das größte Problem.} — À mon avis, le chômage est le plus grand problème.\par
\all{Die Folge ist, dass viele Jugendliche die Schule verlassen.} — La conséquence est que beaucoup de jeunes quittent l'école.
\end{exemplebox}

\begin{center}
\begin{tabularx}{\linewidth}{|X|X|X|}
\hline
\rowcolor{popblueL} Français & Deutsch & Remarque \\
\hline
Le texte traite de... & \all{Der Text handelt von...} & \all{von} + datif \\
\hline
L'auteur montre que... & \all{Der Autor zeigt, dass...} & verbe à la fin après \all{dass} \\
\hline
Je trouve que... & \all{Ich finde, dass...} & verbe à la fin \\
\hline
À mon avis... & \all{Meiner Meinung nach...} & verbe en position 2 \\
\hline
La conséquence est que... & \all{Die Folge ist, dass...} & verbe à la fin \\
\hline
La solution & \all{die Lösung, -en} & féminin \\
\hline
\end{tabularx}
\end{center}

\courssub{Production guidée : un paragraphe argumenté}

Tu vas maintenant rédiger ton propre paragraphe de 5 à 6 phrases sur le thème \all{Ein soziales Problem in meiner Stadt} (Un problème social dans ma ville). La méthode est la même que pour le commentaire, mais appliquée à ta propre réalité : Yaoundé, Douala, Bafoussam...

\begin{methode}[Rédiger un paragraphe argumenté en 5 phrases]
\begin{enumerate}
\item \textbf{Phrase d'introduction} : présente le problème. \all{In meiner Stadt gibt es ein großes Problem: die Arbeitslosigkeit.} (Dans ma ville, il y a un grand problème : le chômage.)
\item \textbf{La cause} : explique pourquoi, avec \all{weil} ou \all{denn}. \all{Viele Jugendliche sind arbeitslos, weil es wenige Fabriken gibt.} (Beaucoup de jeunes sont au chômage parce qu'il y a peu d'usines.)
\item \textbf{La conséquence} : \all{Die Folge ist, dass...} (La conséquence est que...)
\item \textbf{La comparaison} : utilise un comparatif. \all{Das Problem ist heute größer als früher.} (Le problème est aujourd'hui plus grand qu'autrefois.)
\item \textbf{Conclusion personnelle} : ton avis + une solution. \all{Meiner Meinung nach muss die Regierung mehr Schulen bauen.} (À mon avis, le gouvernement doit construire plus d'écoles.)
\end{enumerate}
Consigne de la leçon : utilise \textbf{au moins 2 mots du lexique} (par exemple \all{die Armut}, \all{die Folge}, \all{gefährlich}), \textbf{1 phrase avec \all{weil}}, \textbf{1 comparatif} et \textbf{1 verbe au parfait} (par exemple \all{Die Regierung hat ein Programm gestartet} — le gouvernement a lancé un programme).
\end{methode}

\begin{exoresolu}[Rédiger le paragraphe guidé]
Énoncé : rédige un paragraphe de 5 à 6 phrases « Ein soziales Problem in meiner Stadt » avec 2 mots du lexique, 1 phrase avec \all{weil}, 1 comparatif et 1 verbe au parfait.
\tcblower
\textbf{Solution détaillée (modèle) :}\par
\all{In meiner Stadt Douala gibt es ein soziales Problem: die Armut. Viele Familien sind arm, weil die Eltern keine Arbeit haben. Die Folge ist, dass viele Kinder nicht zur Schule gehen. Das Problem ist heute größer als vor zehn Jahren. Die Regierung hat schon ein Programm gestartet, aber das ist nicht genug. Meiner Meinung nach muss man den Familien mehr helfen.}\par
Traduction : « Dans ma ville Douala, il y a un problème social : la pauvreté. Beaucoup de familles sont pauvres parce que les parents n'ont pas de travail. La conséquence est que beaucoup d'enfants ne vont pas à l'école. Le problème est aujourd'hui plus grand qu'il y a dix ans. Le gouvernement a déjà lancé un programme, mais ce n'est pas suffisant. À mon avis, il faut plus aider les familles. »\par
\textbf{Vérification de la consigne :} mots du lexique : \all{die Armut}, \all{die Folge} (2 mots) ; phrase avec \all{weil} : \all{..., weil die Eltern keine Arbeit haben} (verbe \all{haben} à la fin) ; comparatif : \all{größer als} ; parfait : \all{hat ... gestartet}.
\end{exoresolu}

\begin{attention}[Les trois pièges du paragraphe argumenté]
\begin{itemize}
\item \textbf{Place du verbe} : après \all{weil} et \all{dass}, le verbe va à la fin ; après \all{denn} et \all{Meiner Meinung nach}, il reste en position 2. Vérifie chaque phrase une par une.
\item \textbf{Le parfait} : n'oublie pas le participe à la fin. \all{Die Regierung hat ein Programm \textbf{gestartet}.} — et non \emph{hat gestartet ein Programm}.
\item \textbf{Faux ami} : \all{die Drogen} = la drogue (et non « les médicaments », qui se dit \all{die Medikamente}) ; \all{arm} = pauvre (et non « arme » : \all{die Waffe}).
\end{itemize}
\end{attention}

\courssub{Auto-évaluation et production orale}

Avant de rendre ton paragraphe, contrôle-le toi-même avec cette grille. Coche chaque point : si une réponse est « non », corrige ta phrase.

\begin{center}
\begin{tabularx}{\linewidth}{|X|l|}
\hline
\rowcolor{popblueL} Question d'auto-évaluation & Oui / Non \\
\hline
Ai-je respecté la place du verbe (position 2 ; à la fin après \all{weil}/\all{dass}) ? & \\
\hline
Ai-je utilisé au moins 2 mots du vocabulaire du thème (\all{die Armut}, \all{die Folge}, \all{gefährlich}...) ? & \\
\hline
Ai-je donné une cause avec \all{weil} ou \all{denn} ? & \\
\hline
Ai-je donné une conséquence (\all{Die Folge ist, dass...}) ? & \\
\hline
Ai-je utilisé un comparatif (\all{größer als}, \all{gefährlicher als}) ? & \\
\hline
Ai-je mis un verbe au parfait (auxiliaire + participe à la fin) ? & \\
\hline
Ai-je donné mon avis personnel (\all{Ich finde, dass...} / \all{Meiner Meinung nach...}) ? & \\
\hline
Ai-je mis une majuscule à tous les noms ? & \\
\hline
\end{tabularx}
\end{center}

\begin{experience}[Présentation orale en 1 minute]
Travail en deux temps.\par
\textbf{Étape 1 — Présentation.} Chaque élève se lève et présente son paragraphe à la classe en une minute, sans le lire mot pour mot : il commence par \all{Ich spreche über ein soziales Problem in meiner Stadt: ...} (Je parle d'un problème social dans ma ville : ...), puis donne la cause, la conséquence et son avis.\par
\textbf{Étape 2 — Questions des camarades.} Les autres élèves posent une ou deux questions simples, par exemple : \all{Warum ist das Problem gefährlich?} (Pourquoi le problème est-il dangereux ?), \all{Was kann man machen?} (Qu'est-ce qu'on peut faire ?), \all{Ist das Problem größer als früher?} (Le problème est-il plus grand qu'autrefois ?). L'élève répond avec une phrase complète, en faisant attention à la place du verbe.\par
\textbf{Critères de réussite :} parler fort et lentement, utiliser au moins une phrase-outil de la leçon, répondre aux questions avec un verbe bien placé.
\end{experience}

\begin{aretenir}[L'essentiel de la section]
\begin{itemize}
\item Un commentaire de texte suit 4 étapes : présentation (\all{Der Text handelt von...}), idée principale (\all{Der Autor zeigt, dass...}), analyse (idées secondaires + vocabulaire), avis personnel (\all{Meiner Meinung nach...}).
\item Un paragraphe argumenté suit 5 phrases : problème → cause (\all{weil}/\all{denn}) → conséquence (\all{Die Folge ist, dass...}) → comparaison (\all{größer als...}) → conclusion personnelle.
\item Après \all{dass} et \all{weil}, le verbe conjugué va \textbf{à la fin} ; au parfait, le participe va aussi \textbf{à la fin}.
\item Avant de rendre ta production, vérifie-la avec la grille d'auto-évaluation : place du verbe, vocabulaire du thème, cause, conséquence, comparatif, parfait, avis personnel, majuscules.
\end{itemize}
\end{aretenir}

\newpage

\coursec{Activité d'intégration}

\courssub{Objectifs de l'activité}

À l'issue de cette activité d'intégration, l'élève sera capable de :

\begin{itemize}
\item réinvestir le \cle{vocabulaire thématique} des fléaux sociaux (\all{die Armut, die Drogen, die Gewalt, die Arbeitslosigkeit, die Folge, gefährlich}) dans une production authentique ;
\item exprimer la \cle{cause} avec \all{weil} (subordonnée, verbe à la fin) et \all{denn} (coordonnée, verbe en deuxième position) ;
\item former et employer correctement le \cle{comparatif} (\all{besser als, gefährlicher als, mehr als}) ;
\item utiliser le \cle{parfait} (\all{Perfekt}) avec l'auxiliaire correct (\all{haben} ou \all{sein}) pour raconter un événement passé ;
\item collaborer en groupe, présenter un document à l'oral en deux minutes et justifier un choix en allemand.
\end{itemize}

\courssub{Situation}

L'ONG camerounaise \emph{« Jeunesse Debout »}, en partenariat avec l'association allemande \emph{« Jugend hilft »} de Bonn, organise un concours d'affiches de sensibilisation dans les lycées de Yaoundé et de Douala. Le thème de cette année : \all{„Soziale Probleme — sag NEIN!“} (« Les fléaux sociaux — dis NON ! »). La meilleure affiche sera imprimée et affichée dans les écoles et les maisons des jeunes des deux pays.

Voici la lettre que la directrice de l'ONG a envoyée à votre classe :

\begin{textebox}[Lettre de la directrice de l'ONG]
Liebe Schülerinnen und Schüler,

unsere Welt hat viele soziale Probleme: Armut, Drogen, Gewalt und Arbeitslosigkeit. Viele junge Leute nehmen Drogen, weil sie keine Arbeit haben. Andere sind gewalttätig, denn sie haben keine Perspektive. Das ist gefährlich für die ganze Gesellschaft.

Deshalb organisieren wir einen Wettbewerb. Erstellt ein Plakat gegen ein soziales Problem! Euer Plakat braucht: einen Slogan mit einem Komparativ, zwei Sätze mit „weil“ und „denn“, einen Satz im Perfekt und die Folgen des Problems.

Präsentiert euer Plakat in zwei Minuten auf Deutsch. Die beste Kampagne gewinnt!

Viel Erfolg!

Eure Mama Ngo Bell, Direktorin der NGO „Jeunesse Debout“
\end{textebox}

\textbf{Glossaire :} \all{der Wettbewerb, -e} = le concours ; \all{das Plakat, -e} = l'affiche ; \all{gewalttätig} = violent ; \all{die Perspektive, -n} = la perspective, l'avenir ; \all{die Gesellschaft, -en} = la société ; \all{gewinnen} = gagner ; \all{der Erfolg} = le succès.

\courssub{Consignes}

Travaillez par groupes de quatre élèves. Suivez les étapes dans l'ordre.

\begin{enumerate}
\item \textbf{Choix du fléau (5 min).} Dans votre groupe, choisissez UN fléau social : \all{die Drogen, die Arbeitslosigkeit, die Gewalt} ou \all{die Armut}. Justifiez votre choix à voix haute : \all{Wir wählen die Drogen, weil sie sehr gefährlich sind.}
\item \textbf{Remue-méninges lexical (10 min).} Sur une feuille de brouillon, notez au moins huit mots du thème avec leur article et leur pluriel (ex. \all{die Folge, -n ; der Jugendliche, -n ; die Familie, -n}). Soulignez deux adjectifs utiles (\all{gefährlich, schlimm, wichtig, stark}).
\item \textbf{Le slogan (10 min).} Écrivez un slogan court avec un comparatif et \all{als}. Modèles : \all{Sport ist besser als Drogen!} / \all{Arbeit ist wichtiger als Gewalt!} Attention : le comparatif se forme avec le suffixe \all{-er} (\all{gut} $\rightarrow$ \all{besser}, \all{gefährlich} $\rightarrow$ \all{gefährlicher}).
\item \textbf{Les causes (10 min).} Rédigez deux phrases : l'une avec \all{weil} (le verbe conjugué va à la FIN), l'autre avec \all{denn} (le verbe reste en deuxième position). Exemple : \all{Viele Jugendliche nehmen Drogen, weil sie arbeitslos sind. Sie brauchen Hilfe, denn sie haben keine Perspektive.}
\item \textbf{Le fait passé (5 min).} Ajoutez une phrase au parfait qui décrit ce qui s'est déjà passé : \all{Letztes Jahr haben viele junge Leute ihre Arbeit verloren.} ou \all{Die Gewalt ist in den Städten gestiegen.} Choisissez bien l'auxiliaire : \all{sein} pour les verbes de mouvement ou de changement (\all{steigen, kommen}), \all{haben} pour les autres.
\item \textbf{Les conséquences (10 min).} Écrivez deux ou trois phrases sur les conséquences avec le vocabulaire du thème : \all{Die Folgen sind schlimm: Die Familien werden arm, und die Kinder gehen nicht mehr zur Schule.}
\item \textbf{Fabrication de l'affiche (15 min).} Disposez le slogan en grand, les phrases en dessous, avec des couleurs et un dessin simple. Tout le texte doit être en allemand correct : majuscules aux noms, Umlaute, verbe à la bonne place.
\item \textbf{Présentation orale (2 min par groupe).} Chaque membre du groupe parle au moins 20 secondes. Présentez le fléau, lisez le slogan, expliquez les causes et les conséquences.
\item \textbf{Le vote (10 min).} La classe vote pour la meilleure campagne. Chaque électeur justifie son choix en allemand : \all{Ich finde diese Kampagne besser, weil der Slogan stark ist und die Sätze korrekt sind.}
\end{enumerate}

\courssub{Ressources à mobiliser}

\begin{aretenir}[Boîte à outils de la campagne]
\textbf{Lexique :} \all{die Armut} (pauvreté, sans pluriel), \all{die Droge, -n} (drogue), \all{die Gewalt} (violence), \all{die Arbeitslosigkeit} (chômage), \all{die Folge, -n} (conséquence), \all{gefährlich} (dangereux), \all{der Jugendliche, -n} (jeune), \all{die Hilfe, -n} (aide).

\textbf{La cause :} \all{weil} + verbe à la fin : \all{Er nimmt Drogen, weil er keine Arbeit hat.} ; \all{denn} + verbe en position 2 : \all{Er ist traurig, denn er hat keine Arbeit.}

\textbf{Le comparatif :} adjectif + \all{-er} + \all{als} : \all{Sport ist besser als Drogen.} ; \all{Drogen sind gefährlicher als Alkohol.} ; \all{mehr Geld als}, \all{schlimmer als}.

\textbf{Le parfait :} sujet + \all{haben/sein} + participe II à la fin : \all{Die Jugendlichen haben demonstriert.} / \all{Die Arbeitslosigkeit ist gestiegen.}
\end{aretenir}

\begin{attention}[Pièges à éviter]
\begin{itemize}
\item Après \all{weil}, le verbe conjugué va à la fin : jamais \all{*weil er hat keine Arbeit}, mais \all{weil er keine Arbeit hat}.
\item \all{denn} ne fait pas reculer le verbe : \all{denn er hat keine Arbeit}.
\item Ne confondez pas \all{als} (comparatif : \all{besser als}) et \all{wie} (égalité : \all{so gefährlich wie}).
\item Majuscule à tous les noms : \all{die Drogen, die Gewalt, die Folgen}.
\item Au parfait, \all{steigen} prend \all{sein} : \all{ist gestiegen}.
\end{itemize}
\end{attention}

\courssub{Proposition de production et corrigé commenté}

\begin{exoresolu}[Affiche modèle du groupe 3 : campagne contre la drogue]
Énoncé : rédigez l'affiche complète (slogan, deux causes, un fait passé, les conséquences) et le texte de la présentation orale.
\tcblower
\textbf{L'affiche :}

\all{„Sport ist besser als Drogen!“}

\all{Viele Jugendliche nehmen Drogen, weil sie arbeitslos und frustriert sind. Sie brauchen unsere Hilfe, denn sie haben keine Perspektive. Letztes Jahr ist die Zahl der jungen Drogenkonsumenten in unserer Stadt gestiegen. Die Folgen sind schlimm: Die Familien werden arm, die Gesundheit ist kaputt, und die Kriminalität wird größer. Sag NEIN zu Drogen — sag JA zum Sport!}

\textbf{La présentation orale :}

\all{Guten Tag! Wir präsentieren unsere Kampagne gegen die Drogen. Unser Slogan ist: „Sport ist besser als Drogen!“ Viele Jugendliche nehmen Drogen, weil sie arbeitslos und frustriert sind. Sie brauchen Hilfe, denn sie haben keine Perspektive. Letztes Jahr ist die Zahl der Drogenkonsumenten gestiegen. Die Folgen sind Armut, kaputte Gesundheit und mehr Kriminalität. Deshalb sagen wir: Nein zu Drogen, ja zum Sport! Danke!}

\textbf{Commentaire des choix de langue :}
\begin{itemize}
\item \textbf{Slogan avec comparatif :} \all{besser als} — comparatif irrégulier de \all{gut}, suivi de \all{als} (jamais \all{wie} après un comparatif).
\item \textbf{Cause avec weil :} \all{weil sie arbeitslos und frustriert sind} — le verbe \all{sind} est bien rejeté à la fin de la subordonnée.
\item \textbf{Cause avec denn :} \all{denn sie haben keine Perspektive} — coordination, le verbe \all{haben} reste en deuxième position.
\item \textbf{Parfait :} \all{ist ... gestiegen} — auxiliaire \all{sein} car \all{steigen} exprime un changement ; participe II \all{gestiegen} à la fin.
\item \textbf{Vocabulaire du thème :} \all{die Folgen, die Armut, die Gesundheit, die Kriminalität, die Perspektive} — tous avec majuscule ; \all{größer} est un autre comparatif (de \all{groß}, avec Umlaut).
\item \textbf{Impératif final :} \all{Sag NEIN!} — impératif de la 2\textsuperscript{e} personne du singulier, ton direct et percutant, idéal pour une affiche.
\end{itemize}
\end{exoresolu}

\courssub{Bilan de l'activité}

Cette activité d'intégration a permis de réunir, dans une tâche finale authentique et motivante, tous les savoir-faire de la leçon. D'un point de vue linguistique, les élèves ont mobilisé le lexique des fléaux sociaux avec articles et pluriels, manipulé les deux constructions de la cause (\all{weil} avec verbe final, \all{denn} avec verbe en position 2), formé des comparatifs réguliers et irréguliers (\all{gefährlicher, besser, größer}) et conjugué des verbes au parfait en choisissant le bon auxiliaire. D'un point de vue communicationnel, ils ont travaillé en équipe, négocié des choix, présenté un document à l'oral en respectant un temps imparti et justifié une opinion avec la structure \all{Ich finde ..., weil ...}. Enfin, la dimension citoyenne et interculturelle — une ONG camerounaise associée à une association allemande — montre que l'allemand est un outil réel d'engagement et de dialogue entre le Cameroun et les pays germanophones. Le vote final, argumenté en allemand, clôt la séquence dans une dynamique de classe positive et évaluative.

\newpage

\coursec{Méthodes et exercices résolus}

\courssub{Savoir-faire 1 : Comprendre un texte sur les fléaux sociaux (global puis détaillé)}

\begin{methode}[Lire et comprendre un texte social en trois lectures]
\begin{enumerate}
\item \textbf{Lecture globale} : lire le titre et le premier paragraphe. Identifier le thème (ici : \all{soziale Probleme} — les problèmes sociaux), le type de texte (article, reportage, témoignage) et le sujet principal de chaque paragraphe.
\item \textbf{Lecture détaillée} : souligner les mots-clés du thème (\all{die Armut}, \all{die Arbeitslosigkeit}, \all{die Drogen}, \all{die Gewalt}, \all{die Folge}). Repérer les connecteurs logiques : \all{weil} et \all{denn} (cause), \all{aber} (opposition), \all{deshalb} (conséquence).
\item \textbf{Réponse aux questions} : localiser dans le texte la phrase qui contient la réponse, puis reformuler avec ses propres mots, en phrase complète (sujet + verbe conjugué en 2e position). Ne jamais recopier un long passage.
\item \textbf{Vérification} : relire la réponse : le verbe est-il conjugué et bien placé ? Les noms ont-ils une majuscule ? La réponse répond-elle vraiment à la question (qui ? quoi ? pourquoi ? où ?) ?
\end{enumerate}
\textbf{Réflexes} : une question avec \all{Warum} appelle une réponse avec \all{weil} ou \all{denn} ; une question avec \all{Was} appelle un nom ou une action ; une question avec \all{Wer} appelle une personne.
\end{methode}

\begin{exoresolu}[Compréhension de texte — type Bac]
\textbf{Énoncé.} Lisez l'extrait suivant et répondez en allemand, par des phrases complètes.

\begin{textebox}[Soziale Probleme in unserer Welt (extrait)]
Viele Jugendliche in den großen Städten haben keine Arbeit. Weil sie arbeitslos sind, haben sie oft kein Geld für die Schule oder für das Essen. Einige nehmen dann Drogen, denn sie sehen keine Zukunft mehr. Die Folgen sind gefährlich: Gewalt in den Familien, Kriminalität in den Straßen und Krankheiten. Die Regierung und die Familien müssen zusammenarbeiten, weil nur so eine Lösung möglich ist.
\end{textebox}

\textbf{Questions :}
\begin{enumerate}
\item Warum haben viele Jugendliche kein Geld für die Schule ?
\item Warum nehmen einige Jugendliche Drogen ?
\item Was sind die Folgen dieser Probleme ? (Nennen Sie zwei Folgen.)
\end{enumerate}
\tcblower
\textbf{Solution détaillée.}

\textbf{Question 1.} La question commence par \all{Warum} : la réponse exige une cause. Dans le texte : \all{Weil sie arbeitslos sind, haben sie oft kein Geld...}. On reformule en phrase complète avec \all{weil} (verbe en fin de proposition subordonnée) :

\all{Viele Jugendliche haben kein Geld für die Schule, weil sie arbeitslos sind.}

« Beaucoup de jeunes n'ont pas d'argent pour l'école parce qu'ils sont au chômage. »

Justification : \all{weil} + sujet \all{sie} + verbe conjugué \all{sind} rejeté à la fin. On aurait pu utiliser \all{denn} : \all{..., denn sie sind arbeitslos} (verbe en 2e position, car \all{denn} est une coordination).

\textbf{Question 2.} Le texte dit : \all{..., denn sie sehen keine Zukunft mehr}. Reformulation :

\all{Einige Jugendliche nehmen Drogen, weil sie keine Zukunft mehr sehen.}

« Certains jeunes prennent de la drogue parce qu'ils ne voient plus d'avenir. »

Attention à la place de \all{mehr} : il se place avant le participe/infinitif en fin de phrase (\all{nicht mehr sehen}).

\textbf{Question 3.} La question avec \all{Was} demande une énumération. Le texte donne trois folgen ; on en choisit deux :

\all{Die Folgen sind Gewalt in den Familien und Kriminalität in den Straßen.}

« Les conséquences sont la violence dans les familles et la criminalité dans les rues. »

\textbf{Conclusion méthodologique} : chaque réponse est une phrase complète, avec verbe conjugué en 2e position, majuscules aux noms (\all{Geld}, \all{Schule}, \all{Zukunft}, \all{Folgen}, \all{Gewalt}) et connecteur de cause correctement employé.
\end{exoresolu}

\courssub{Savoir-faire 2 : Employer le vocabulaire thématique des fléaux sociaux}

\begin{methode}[Construire et utiliser le champ lexical des problèmes sociaux]
\begin{enumerate}
\item \textbf{Apprendre chaque nom avec son article et son pluriel} : \all{die Armut} (–, pauvreté), \all{die Arbeitslosigkeit} (–, chômage), \all{die Gewalt} (–, violence), \all{die Droge, -n} (drogue), \all{die Folge, -n} (conséquence), \all{das Problem, -e} (problème), \all{die Lösung, -en} (solution).
\item \textbf{Mémoriser les familles de mots} : \all{die Arbeit} → \all{arbeitslos} (adj.) → \all{die Arbeitslosigkeit} ; \all{die Gefahr} → \all{gefährlich} (dangereux) ; \all{arm} → \all{die Armut}.
\item \textbf{Associer nom + verbe} (collocations) : \all{Drogen nehmen} (prendre de la drogue), \all{ein Problem lösen} (résoudre un problème), \all{Arbeit suchen / finden} (chercher / trouver du travail), \all{Folgen haben} (avoir des conséquences).
\item \textbf{Réutiliser le lexique dans ses propres phrases} : pour chaque mot appris, écrire une phrase personnelle sur le contexte camerounais (Yaoundé, Douala, les écoles).
\end{enumerate}
\textbf{Réflexe} : au Bac, un mot sans article ou avec un mauvais genre coûte des points. Toujours vérifier : \all{die Armut}, jamais \all{der Armut}.
\end{methode}

\begin{exoresolu}[Lexique — compléter un texte à trous]
\textbf{Énoncé.} Complétez le texte avec les mots suivants (avec le bon article si nécessaire) : \all{Arbeitslosigkeit}, \all{gefährlich}, \all{Drogen}, \all{Folgen}, \all{Lösung}.

\all{In Douala ist die (1) \dots\ ein großes Problem. Viele junge Leute finden keine Arbeit. Manche nehmen (2) \dots\ , und das ist sehr (3) \dots\ für die Gesundheit. Die (4) \dots\ sind schlecht für die ganze Familie. Die Stadt sucht eine (5) \dots\ für dieses Problem.}
\tcblower
\textbf{Solution détaillée.}
\begin{enumerate}
\item \all{die Arbeitslosigkeit} : le contexte « viele junge Leute finden keine Arbeit » (beaucoup de jeunes ne trouvent pas de travail) désigne le chômage. Article féminin déjà donné (\all{die}), nom avec majuscule.
\item \all{Drogen} : verbe \all{nehmen} + complément pluriel sans article (sens général : « prendre de la drogue »). Majuscule obligatoire.
\item \all{gefährlich} : adjectif prédicatif après \all{ist}, donc \textbf{invariable} (pas de terminaison). « cela est très dangereux pour la santé ».
\item \all{Die Folgen} : pluriel de \all{die Folge} ; le verbe \all{sind} confirme le pluriel. « Les conséquences sont mauvaises pour toute la famille. »
\item \all{eine Lösung} : après l'article indéfini \all{eine} (féminin, accusatif), le nom \all{Lösung} convient : « La ville cherche une solution à ce problème. »
\end{enumerate}
\textbf{Texte corrigé complet :} \all{In Douala ist die Arbeitslosigkeit ein großes Problem. Viele junge Leute finden keine Arbeit. Manche nehmen Drogen, und das ist sehr gefährlich für die Gesundheit. Die Folgen sind schlecht für die ganze Familie. Die Stadt sucht eine Lösung für dieses Problem.}

\textbf{Conclusion} : le choix du mot dépend du sens (contexte) et de la grammaire (article, pluriel, adjectif prédicatif invariable).
\end{exoresolu}

\courssub{Savoir-faire 3 : Exprimer la cause avec weil et denn (thème grammatical)}

\begin{methode}[Choisir et construire une phrase de cause]
\begin{enumerate}
\item \textbf{Identifier la relation logique} : la 2e phrase explique la 1re (question : « pourquoi ? »).
\item \textbf{Choisir le connecteur} :
\begin{itemize}
\item \all{weil} = subordination → le verbe conjugué va \textbf{à la fin} : \all{Er bleibt zu Hause, weil er krank ist.}
\item \all{denn} = coordination → le verbe conjugué reste en \textbf{2e position} : \all{Er bleibt zu Hause, denn er ist krank.}
\end{itemize}
\item \textbf{Traduire « parce que / car »} : les deux se traduisent par « parce que » ; \all{denn} est un peu plus soutenu (« car »).
\item \textbf{Vérifier} : après \all{weil}, compter les éléments — le verbe conjugué doit être le dernier mot de la subordonnée ; si la subordonnée ouvre la phrase, le verbe de la principale vient juste après la virgule : \all{Weil er krank ist, bleibt er zu Hause.}
\end{enumerate}
\textbf{Structure à retenir} : \all{Weil + S + ... + V, V + S + ...} (subordonnée en tête = le verbe de la principale occupe la 1re position après la virgule).
\end{methode}

\begin{exoresolu}[Thème — traduction de phrases avec la cause]
\textbf{Énoncé.} Traduisez en allemand.
\begin{enumerate}
\item Beaucoup d'enfants ne vont pas à l'école parce que leurs parents sont pauvres. (avec \all{weil})
\item Paul ne prend pas de drogue, car il connaît les conséquences dangereuses. (avec \all{denn})
\item Parce que le chômage est un grand problème, le gouvernement cherche des solutions. (\all{weil} en tête de phrase)
\end{enumerate}
\tcblower
\textbf{Solution détaillée.}

\textbf{1.} \all{Viele Kinder gehen nicht in die Schule, weil ihre Eltern arm sind.}

Raisonnement : « parce que » = \all{weil} → subordonnée, verbe \all{sind} à la fin. « leurs parents » = \all{ihre Eltern} (possessif pluriel). « pauvres » = \all{arm} (adjectif prédicatif invariable). « aller à l'école » = \all{in die Schule gehen} (accusatif de direction).

\textbf{2.} \all{Paul nimmt keine Drogen, denn er kennt die gefährlichen Folgen.}

Raisonnement : « car » = \all{denn} → coordination, verbe \all{kennt} en 2e position. « ne pas prendre de drogue » = négation d'un complément sans article : \all{keine Drogen nehmen}. « les conséquences dangereuses » = \all{die gefährlichen Folgen} : adjectif épithète à l'accusatif pluriel après article défini → terminaison \all{-en}.

\textbf{3.} \all{Weil die Arbeitslosigkeit ein großes Problem ist, sucht die Regierung Lösungen.}

Raisonnement : la subordonnée \all{weil} ouvre la phrase → verbe \all{ist} à la fin de la subordonnée, puis virgule, puis le verbe de la principale \all{sucht} immédiatement en 1re position (le sujet \all{die Regierung} suit le verbe : inversion obligatoire). « le chômage » = \all{die Arbeitslosigkeit} ; « un grand problème » = \all{ein großes Problem} (neutre, adjectif en \all{-es} après \all{ein}).

\textbf{Conclusion} : le choix \all{weil}/\all{denn} détermine toute la mécanique de la phrase : fin de phrase pour \all{weil}, 2e position pour \all{denn}.
\end{exoresolu}

\courssub{Savoir-faire 4 : Former et employer le comparatif}

\begin{methode}[Construire une comparaison en allemand]
\begin{enumerate}
\item \textbf{Former le comparatif} : adjectif + \all{-er} : \all{arm} → \all{ärmer} ; \all{gefährlich} → \all{gefährlicher} ; \all{groß} → \all{größer}. Les adjectifs courts à voyelle \all{a, o, u} prennent souvent un tréma : \all{alt} → \all{älter}, \all{jung} → \all{jünger}, \all{arm} → \all{ärmer}. Irréguliers : \all{gut} → \all{besser}, \all{viel} → \all{mehr}, \all{gern} → \all{lieber}.
\item \textbf{Comparer deux éléments} : \all{als} (= « que ») après le comparatif : \all{Die Stadt ist gefährlicher als das Dorf.} « La ville est plus dangereuse que le village. »
\item \textbf{Égalité} : \all{so ... wie} : \all{Paul ist so arm wie sein Bruder.} « Paul est aussi pauvre que son frère. »
\item \textbf{Accorder si l'adjectif est épithète} : \all{ein gefährlicheres Problem} (un problème plus dangereux) : le comparatif prend les désinences normales de l'adjectif.
\end{enumerate}
\textbf{Réflexe} : « que » de comparaison = \all{als}, jamais \all{wie} ; « aussi ... que » = \all{so ... wie}, jamais \all{als}.
\end{methode}

\begin{exoresolu}[Comparatif — transformation et traduction]
\textbf{Énoncé A.} Transformez selon le modèle : \all{Die Stadt ist gefährlich. Das Dorf ist nicht so gefährlich.} → \all{Die Stadt ist gefährlicher als das Dorf.}
\begin{enumerate}
\item \all{Die Arbeitslosigkeit ist groß. Die Armut ist nicht so groß.} (comparer avec \all{groß})
\item \all{Drogen sind schlecht für die Gesundheit. Sport ist gut.} (comparer avec \all{gut})
\end{enumerate}
\textbf{Énoncé B.} Traduisez : « La vie à Yaoundé est plus chère que la vie au village, mais les jeunes y trouvent plus de travail. »
\tcblower
\textbf{Solution détaillée.}

\textbf{A.1.} \all{groß} → comparatif \all{größer} (tréma sur \all{o}, adjectif monosyllabique). Phrase : \all{Die Arbeitslosigkeit ist größer als die Armut.} « Le chômage est plus important que la pauvreté. » Adjectif prédicatif → \all{größer} reste invariable.

\textbf{A.2.} \all{gut} est irrégulier : \all{besser}. On compare \all{Sport} à \all{Drogen} : \all{Sport ist besser für die Gesundheit als Drogen.} « Le sport est meilleur pour la santé que la drogue. »

\textbf{B.} \all{Das Leben in Yaoundé ist teurer als das Leben auf dem Dorf, aber die Jugendlichen finden dort mehr Arbeit.}

Raisonnement : « la vie » = \all{das Leben} (neutre) ; « plus chère » = \all{teurer} (comparatif de \all{teuer}, on perd le \all{e} : \all{teuer} → \all{teurer}) ; « que » = \all{als} ; « au village » = \all{auf dem Dorf} (datif de lieu) ; « y » = \all{dort} ; « plus de travail » = \all{mehr Arbeit} (comparatif irrégulier de \all{viel}). Coordination avec \all{aber} : verbe \all{finden} en 2e position.

\textbf{Conclusion} : trois réflexes — tréma sur les voyelles fortes, irréguliers \all{besser}/\all{mehr} par cœur, \all{als} pour l'inégalité.
\end{exoresolu}

\courssub{Savoir-faire 5 : Raconter au parfait (Perfekt)}

\begin{methode}[Construire le Perfekt pour raconter des faits passés]
\begin{enumerate}
\item \textbf{Choisir l'auxiliaire} : \all{haben} pour la plupart des verbes ; \all{sein} pour les verbes de déplacement ou de changement d'état (\all{gehen}, \all{kommen}, \all{sterben}, \all{aufstehen}) et pour \all{sein}/\all{bleiben}.
\item \textbf{Former le participe II} : verbes faibles : ge + radical + t (\all{machen} → \all{gemacht}) ; verbes forts : ge + (radical modifié) + en (\all{nehmen} → \all{genommen}, \all{finden} → \all{gefunden}, \all{helfen} → \all{geholfen}) ; verbes en \all{-ieren} : pas de \all{ge-} (\all{studieren} → \all{studiert}).
\item \textbf{Placer les éléments} : auxiliaire conjugué en 2e position, participe II \textbf{à la fin} : \all{Er hat Drogen genommen.} / \all{Sie ist in die Stadt gegangen.}
\item \textbf{Verbes à particule séparable} : \all{ge-} s'insère entre la particule et le radical : \all{aufstehen} → \all{aufgestanden}.
\end{enumerate}
\textbf{Réflexe} : au passé composé du récit oral/écrit simple, le Perfekt est le temps par défaut ; réserver le Präteritum à \all{sein} (\all{war}), \all{haben} (\all{hatte}) et aux modaux.
\end{methode}

\begin{exoresolu}[Conjugaison — récit au Perfekt et version]
\textbf{Énoncé A.} Mettez au Perfekt :
\begin{enumerate}
\item \all{Der Junge nimmt Drogen.}
\item \all{Die Familie zieht in die Stadt.} (déménager : \all{ziehen}, verbe fort, auxiliaire \all{sein})
\item \all{Die Regierung hilft den armen Familien.}
\end{enumerate}
\textbf{Énoncé B.} Traduisez en français : \all{Viele Jugendliche sind nach Douala gekommen, weil sie dort Arbeit gesucht haben. Aber sie haben keine Arbeit gefunden, und einige haben leider Drogen genommen.}
\tcblower
\textbf{Solution détaillée.}

\textbf{A.1.} \all{nehmen} est un verbe fort (\all{nimmt} — \all{nahm} — \all{genommen}), auxiliaire \all{haben} : \all{Der Junge hat Drogen genommen.} « Le garçon a pris de la drogue. » Auxiliaire en 2e position, participe \all{genommen} en fin.

\textbf{A.2.} \all{ziehen} (déménager) exprime un déplacement → auxiliaire \all{sein} ; participe fort \all{gezogen} : \all{Die Familie ist in die Stadt gezogen.} « La famille a déménagé en ville. »

\textbf{A.3.} \all{helfen} est fort (\all{hilft} — \all{half} — \all{geholfen}), auxiliaire \all{haben} ; \all{helfen} exige le \textbf{datif} : \all{den armen Familien} (datif pluriel en \all{-n}) : \all{Die Regierung hat den armen Familien geholfen.} « Le gouvernement a aidé les familles pauvres. »

\textbf{B.} Traduction mot à mot contrôlée : \all{sind ... gekommen} = « sont venus » (auxiliaire \all{sein} car déplacement) ; \all{weil sie dort Arbeit gesucht haben} = « parce qu'ils y ont cherché du travail » (subordonnée : auxiliaire \all{haben} en toute fin, après le participe) ; \all{haben keine Arbeit gefunden} = « n'ont pas trouvé de travail » ; \all{haben leider Drogen genommen} = « ont malheureusement pris de la drogue ».

Traduction finale : « Beaucoup de jeunes sont venus à Douala parce qu'ils y cherchaient du travail. Mais ils n'ont pas trouvé de travail, et certains ont malheureusement pris de la drogue. »

\textbf{Conclusion} : trois contrôles systématiques — bon auxiliaire, participe correct (fort/faible), place : auxiliaire en 2e position et participe en fin (auxiliaire après le participe dans une subordonnée \all{weil}).
\end{exoresolu}

\courssub{Savoir-faire 6 : Produire un court paragraphe (résumé et opinion)}

\begin{methode}[Rédiger un court texte sur un fléau social en 5 étapes]
\begin{enumerate}
\item \textbf{Plan simple en 3 parties} : présenter le problème (1 phrase) → expliquer les causes (1-2 phrases avec \all{weil}/\all{denn}) → proposer une solution ou donner son avis (1-2 phrases avec \all{meiner Meinung nach}, \all{ich denke, dass ...}).
\item \textbf{Recycler la grammaire de la leçon} : au moins une phrase avec \all{weil}, une avec \all{denn}, un comparatif, un verbe au Perfekt.
\item \textbf{Utiliser le lexique appris} : \all{die Armut}, \all{die Arbeitslosigkeit}, \all{die Drogen}, \all{die Gewalt}, \all{die Folgen}, \all{gefährlich}, \all{die Lösung}.
\item \textbf{Rédiger} des phrases courtes (8-12 mots) : verbe en 2e position, majuscules aux noms.
\item \textbf{Relire} avec la grille : verbe conjugué ? place du verbe ? articles et majuscules ? participe II en fin au Perfekt ?
\end{enumerate}
\textbf{Phrases-modèles} : \all{Ein großes Problem in Kamerun ist ...} / \all{Viele Jugendliche ..., weil ...} / \all{Meiner Meinung nach muss die Regierung ...}
\end{methode}

\begin{exoresolu}[Production écrite guidée — type Bac]
\textbf{Énoncé.} Rédigez un court paragraphe (5-6 phrases, 50-70 mots) sur le sujet : \all{Die Arbeitslosigkeit in meinem Land}. Consignes : utilisez au moins une phrase avec \all{weil}, une avec \all{denn}, un comparatif et un verbe au Perfekt.
\tcblower
\textbf{Solution détaillée — production modèle.}

\all{Ein großes Problem in Kamerun ist die Arbeitslosigkeit. Viele Jugendliche haben die Schule beendet, aber sie finden keine Arbeit. Sie sind oft frustriert, weil sie kein Geld verdienen. Manche gehen in die großen Städte, denn dort ist das Leben besser als im Dorf. Meiner Meinung nach muss die Regierung mehr Jobs schaffen und den Jugendlichen helfen.}

« Un grand problème au Cameroun est le chômage. Beaucoup de jeunes ont terminé l'école, mais ils ne trouvent pas de travail. Ils sont souvent frustrés parce qu'ils ne gagnent pas d'argent. Certains partent dans les grandes villes, car la vie y est meilleure qu'au village. À mon avis, le gouvernement doit créer plus d'emplois et aider les jeunes. »

\textbf{Vérification des consignes} :
\begin{itemize}
\item \all{weil} : \all{..., weil sie kein Geld verdienen} — verbe \all{verdienen} en fin. \checkmark
\item \all{denn} : \all{..., denn dort ist das Leben besser} — verbe \all{ist} en 2e position. \checkmark
\item Comparatif : \all{besser als im Dorf} (irrégulier \all{gut} → \all{besser} ; \all{im} = \all{in dem}, datif). \checkmark
\item Perfekt : \all{haben die Schule beendet} — auxiliaire \all{haben} en 2e position, participe \all{beendet} en fin (verbe en \all{-ieren} ? Non : \all{beenden} est faible avec préfixe inséparable \all{be-}, donc participe sans \all{ge-} : \all{beendet}). \checkmark
\end{itemize}
\textbf{Conclusion} : un paragraphe court, structuré, qui recycle exactement les points de grammaire et le lexique de la leçon, obtient la totalité des points de production écrite.
\end{exoresolu}

\begin{attention}[Les 5 erreurs qui coûtent des points au Bac]
\begin{enumerate}
\item \textbf{Oublier la majuscule aux noms} : écrire \all{die armut} au lieu de \all{die Armut}. En allemand, \textbf{tous} les noms prennent une majuscule : \all{das Problem}, \all{die Drogen}, \all{die Gewalt}. C'est la faute la plus pénalisée.
\item \textbf{Mal placer le verbe après weil} : \all{..., weil sie sind arbeitslos} est faux. Avec \all{weil}, le verbe conjugué va \textbf{à la fin} : \all{..., weil sie arbeitslos sind.} Ne pas confondre avec \all{denn} (verbe en 2e position).
\item \textbf{Faux amis et calques du français} : \all{die Drogen} signifie « la drogue » (stupéfiants), pas « les médicaments » (\all{die Medikamente}) ; « actuellement » ne se traduit pas par \all{aktuell} dans tous les contextes ; « aider quelqu'un » = \all{jemandem helfen} (\textbf{datif}), jamais l'accusatif : \all{Die Regierung hilft den Familien.}
\item \textbf{Confondre als et wie dans la comparaison} : « plus dangereux que » = \all{gefährlicher als} ; « aussi dangereux que » = \all{so gefährlich wie}. Écrire \all{gefährlicher wie} est une faute grave.
\item \textbf{Perfekt mal construit} : mauvais auxiliaire (\all{Er hat gekommen} au lieu de \all{Er ist gekommen} — verbe de déplacement → \all{sein}), participe oublié ou mal placé (\all{Er hat genommen Drogen} au lieu de \all{Er hat Drogen genommen}), ou \all{ge-} ajouté à tort aux verbes en \all{-ieren} et aux préfixes inséparables (\all{studiert}, \all{beendet}, jamais \all{gestudiert}/\all{gebeendet}).
\end{enumerate}
\end{attention}

\newpage

\coursec{Exercices}

\courssub{Je vérifie mes connaissances}

\begin{exercice}[QCM : le bon choix]
Entoure la bonne réponse.
\begin{enumerate}
\item \all{Die Armut} bedeutet : a) la violence \quad b) la pauvreté \quad c) la maladie
\item \all{Er bleibt zu Hause, \ldots\ er krank ist.} a) \all{denn} \quad b) \all{weil} \quad c) \all{und}
\item Le comparatif de \all{gut} est : a) \all{guter} \quad b) \all{besser} \quad c) \all{mehr gut}
\item \all{Ich \ldots\ gestern ins Krankenhaus gegangen.} a) \all{habe} \quad b) \all{sein} \quad c) \all{bin}
\end{enumerate}
\end{exercice}

\begin{exercice}[Vrai ou faux ?]
Réponds par \all{richtig} ou \all{falsch} et justifie ta réponse par une phrase en allemand.
\begin{enumerate}
\item Après \all{denn}, le verbe conjugué se place à la fin de la phrase.
\item Après \all{weil}, le verbe conjugué se place à la fin de la proposition.
\item Le Partizip II de \all{kommen} est \all{gekommt}.
\item \all{gefährlich} signifie « généreux ».
\end{enumerate}
\end{exercice}

\begin{exercice}[Vocabulaire : les fléaux sociaux]
Complète chaque phrase avec le mot qui convient : \all{die Armut, die Drogen, die Gewalt, die Arbeitslosigkeit, die Folge, gefährlich}.
\begin{enumerate}
\item Viele Jugendliche in Douala haben keinen Job : \all{die \ldots} ist ein großes Problem.
\item \all{\ldots} machen krank und können tödlich sein.
\item Ohne Geld kann man nichts kaufen : das ist \all{die \ldots}.
\item Kämpfe in der Schule sind \all{\ldots} : das ist \all{die \ldots}.
\end{enumerate}
\end{exercice}

\begin{exercice}[Conjugaison : le parfait]
Complète avec l'auxiliaire \all{haben} ou \all{sein} et le Partizip II du verbe entre parenthèses.
\begin{enumerate}
\item Wir \ldots\ einen Text über soziale Probleme \ldots\ (lesen).
\item Die Familie \ldots\ nach Yaoundé \ldots\ (ziehen).
\item Er \ldots\ viel Geld \ldots\ (verlieren).
\item Die Jugendlichen \ldots\ zu Hause \ldots\ (bleiben).
\end{enumerate}
\end{exercice}

\courssub{Je m'entraîne}

\begin{exercice}[Traduction : thème]
Traduis en allemand.
\begin{enumerate}
\item Beaucoup de jeunes sont pauvres parce qu'ils n'ont pas de travail.
\item La drogue est plus dangereuse que l'alcool.
\item Hier, j'ai vu un reportage sur la violence dans les écoles.
\item Le chômage a de mauvaises conséquences pour les familles.
\end{enumerate}
\end{exercice}

\begin{exercice}[Transformation : \all{denn} et \all{weil}]
Réécris chaque phrase en remplaçant \all{denn} par \all{weil} (et inversement). Attention à la place du verbe !
\begin{enumerate}
\item \all{Er lernt Deutsch, denn er will in Deutschland studieren.}
\item \all{Sie bleibt zu Hause, weil sie krank ist.}
\item \all{Wir helfen den Armen, denn sie haben kein Geld.}
\item \all{Die Kinder gehen nicht zur Schule, weil die Familie arm ist.}
\end{enumerate}
\end{exercice}

\begin{exercice}[Texte à trous : le comparatif]
Complète le texte avec le comparatif des adjectifs entre parenthèses.

\all{In der Stadt ist das Leben \ldots\ (teuer) als auf dem Land. Die Häuser sind \ldots\ (groß), aber viele Menschen sind \ldots\ (arm). Die Straßen sind nachts \ldots\ (gefährlich). Auf dem Land ist die Luft \ldots\ (gut) und die Menschen sind oft \ldots\ (freundlich).}
\end{exercice}

\begin{exercice}[Appariement : mots et définitions]
Relie chaque mot allemand à sa définition française.
\begin{enumerate}
\item \all{die Arbeitslosigkeit} \hspace{1cm} a) la conséquence
\item \all{die Folge} \hspace{1cm} b) le chômage
\item \all{die Gewalt} \hspace{1cm} c) dangereux
\item \all{gefährlich} \hspace{1cm} d) la violence
\item \all{die Armut} \hspace{1cm} e) la pauvreté
\end{enumerate}
\end{exercice}

\courssub{Je me prépare au Bac}

\begin{exercice}[Type Bac — Compréhension écrite]
Lis le texte suivant, puis réponds aux questions.

\begin{textebox}[Jugendarbeitslosigkeit in Deutschland]
\all{Viele junge Menschen in Deutschland suchen nach der Schule eine Arbeit. Das ist nicht immer leicht, denn die Firmen wollen oft Arbeiter mit Erfahrung. Die Jugendarbeitslosigkeit ist ein soziales Problem, weil junge Menschen ohne Arbeit kein Geld verdienen können. Sie bleiben länger bei den Eltern, weil eine eigene Wohnung zu teuer ist. Manche Jugendliche sind traurig, andere werden krank. Die Folgen sind gefährlich : Armut, Stress und manchmal auch Drogen. Die Regierung hat deshalb Programme gestartet. Jugendliche lernen einen Beruf in einer Firma und in der Schule. Diese Ausbildung ist besser als nichts, denn die Jugendlichen bekommen Erfahrung und ein Diplom. Viele finden danach eine gute Arbeit und verdienen mehr Geld als vorher. Experten sagen : Bildung ist der beste Weg gegen die Arbeitslosigkeit.}
\end{textebox}

\textbf{A. Questions de compréhension} — Réponds en allemand par des phrases complètes.
\begin{enumerate}
\item \all{Warum ist es für junge Menschen schwer, eine Arbeit zu finden?}
\item \all{Warum bleiben viele Jugendliche bei den Eltern?}
\item \all{Welche Folgen hat die Arbeitslosigkeit? Nenne zwei.}
\item \all{Was lernen die Jugendlichen in der Ausbildung?}
\item \all{Was sagen die Experten über Bildung?}
\end{enumerate}

\textbf{B. Vrai ou faux ?} — Réponds par \all{richtig} ou \all{falsch} et justifie avec une phrase du texte.
\begin{enumerate}
\item \all{Die Firmen wollen Arbeiter ohne Erfahrung.}
\item \all{Eine eigene Wohnung ist billig.}
\item \all{Nach der Ausbildung finden viele Jugendliche eine Arbeit.}
\end{enumerate}

\textbf{C. Questions de langue}
\begin{enumerate}
\item Trouve dans le texte une subordonnée avec \all{weil} et recopie-la.
\item Trouve dans le texte deux comparatifs et donne l'adjectif de base.
\end{enumerate}
\end{exercice}

\begin{exercice}[Type Bac — Grammaire et conjugaison]
\textbf{A. La cause.} Complète les six phrases avec \all{weil} ou \all{denn} et place correctement le verbe conjugué.
\begin{enumerate}
\item \all{Die Kinder hungern, \ldots\ die Familie kein Geld (haben).}
\item \all{Er nimmt keine Drogen, \ldots\ sie gefährlich (sein).}
\item \all{Viele Jugendliche sind arbeitslos, \ldots\ es wenige Jobs (geben).}
\item \all{Wir bleiben zu Hause, \ldots\ es regnet.}
\item \all{Sie lernt viel, \ldots\ sie eine gute Note (wollen).}
\item \all{Die Stadt ist laut, \ldots\ es viele Autos (haben).}
\end{enumerate}

\textbf{B. Le parfait.} Mets les phrases au parfait : choisis le bon auxiliaire (\all{haben} ou \all{sein}) et le bon Partizip II.
\begin{enumerate}
\item \all{Der Junge (finden) keine Arbeit.}
\item \all{Die Familie (ziehen) in die Stadt.}
\item \all{Wir (helfen) den armen Kindern.}
\item \all{Sie (bleiben) lange im Krankenhaus.}
\item \all{Ich (sehen) einen Film über Drogen.}
\item \all{Die Schüler (fahren) nach Deutschland.}
\end{enumerate}

\textbf{C. Version.} Traduis en français l'extrait suivant :

\all{„Die Jugendarbeitslosigkeit ist ein soziales Problem, weil junge Menschen ohne Arbeit kein Geld verdienen können. Die Folgen sind gefährlich : Armut, Stress und manchmal auch Drogen. Diese Ausbildung ist besser als nichts, denn die Jugendlichen bekommen Erfahrung und ein Diplom.“}
\end{exercice}

\begin{exercice}[Type Bac — Expression écrite guidée]
Rédige en allemand un paragraphe de 5 à 6 phrases (environ 50 à 60 mots) sur le sujet suivant :

\textbf{\all{Ein soziales Problem in Kamerun}}

Consignes obligatoires :
\begin{itemize}
\item utilise au moins \textbf{une phrase avec \all{weil}} ;
\item utilise au moins \textbf{un comparatif} ;
\item utilise au moins \textbf{un verbe au parfait} ;
\item utilise au moins \textbf{trois mots du lexique} de la leçon (\all{die Armut, die Drogen, die Gewalt, die Arbeitslosigkeit, die Folge, gefährlich}) ;
\item termine par une phrase qui propose une solution.
\end{itemize}
\end{exercice}

\newpage

\coursec{Corrigés des exercices}

\begin{corrige}[Exercice 1 — QCM : le bon choix]
\begin{enumerate}
\item \textbf{b)} \all{Die Armut} = la pauvreté. (a = \all{die Gewalt}, c = \all{die Krankheit}.)
\item \textbf{a)} \all{Er bleibt zu Hause, denn er krank ist.} — Après \all{denn} (coordinante), le verbe conjugué reste en position 2, comme ici \all{ist}. Avec \all{weil}, il faudrait : \all{\ldots, weil er krank ist}.
\item \textbf{b)} Le comparatif de \all{gut} est irrégulier : \all{gut — besser — am besten}.
\item \textbf{c)} \all{Ich bin gestern ins Krankenhaus gegangen.} — \all{gehen} exprime un déplacement : auxiliaire \all{sein} au parfait.
\end{enumerate}
\end{corrige}

\begin{corrige}[Exercice 2 — Vrai ou faux ?]
\begin{enumerate}
\item \textbf{falsch} — \all{Nach „denn“ steht das Verb an zweiter Stelle : „Er bleibt zu Hause, denn er ist krank.“} (Après \all{denn}, le verbe est en position 2 : c'est une coordination.)
\item \textbf{richtig} — \all{Nach „weil“ steht das Verb am Ende : „\ldots, weil er krank ist.“} (\all{weil} introduit une subordonnée : verbe conjugué rejeté à la fin.)
\item \textbf{richtig} — \all{Das Partizip II von „kommen“ ist „gekommen“ : „Er ist nach Yaoundé gekommen.“} (verbe fort : \all{kommen — kam — ist gekommen}.)
\item \textbf{falsch} — \all{„Gefährlich“ bedeutet „dangereux“, nicht „généreux“.} (Attention au faux ami : « généreux » se dit \all{großzügig}.)
\end{enumerate}
\end{corrige}

\begin{corrige}[Exercice 3 — Vocabulaire : les fléaux sociaux]
\begin{enumerate}
\item \all{Viele Jugendliche in Douala haben keinen Job : die \textbf{Arbeitslosigkeit} ist ein großes Problem.} (le chômage)
\item \all{\textbf{Drogen} machen krank und können tödlich sein.} (les drogues ; pluriel de \all{die Droge})
\item \all{Ohne Geld kann man nichts kaufen : das ist die \textbf{Armut}.} (la pauvreté)
\item \all{Kämpfe in der Schule sind \textbf{gefährlich} : das ist die \textbf{Gewalt}.} (dangereux ; la violence)
\end{enumerate}
\textbf{Points de vigilance :} majuscule à tous les noms (\all{Arbeitslosigkeit, Drogen, Armut, Gewalt}) ; \all{gefährlich} est un adjectif, donc sans majuscule.
\end{corrige}

\begin{corrige}[Exercice 4 — Conjugaison : le parfait]
\begin{enumerate}
\item \all{Wir \textbf{haben} einen Text über soziale Probleme \textbf{gelesen}.} — \all{lesen} : verbe sans déplacement ni changement d'état $\Rightarrow$ auxiliaire \all{haben} ; Partizip II irrégulier \all{gelesen}.
\item \all{Die Familie \textbf{ist} nach Yaoundé \textbf{gezogen}.} — \all{ziehen} exprime ici un déplacement (déménager) $\Rightarrow$ auxiliaire \all{sein}.
\item \all{Er \textbf{hat} viel Geld \textbf{verloren}.} — \all{verlieren} : verbe fort avec préfixe inséparable \all{ver-} $\Rightarrow$ Partizip II sans \all{ge-} : \all{verloren} ; auxiliaire \all{haben}.
\item \all{Die Jugendlichen \textbf{sind} zu Hause \textbf{geblieben}.} — \all{bleiben} prend toujours \all{sein} au parfait : \all{bleiben — blieb — ist geblieben}.
\end{enumerate}
\end{corrige}

\begin{corrige}[Exercice 5 — Traduction : thème]
\begin{enumerate}
\item \all{Viele Jugendliche sind arm, weil sie keine Arbeit haben.} — Verbe \all{haben} rejeté à la fin après \all{weil} ; « pas de travail » = \all{keine Arbeit} (accusatif féminin).
\item \all{Die Droge ist gefährlicher als der Alkohol.} — Comparatif : \all{gefährlich + -er} ; la comparaison se fait avec \all{als} (jamais \all{wie}).
\item \all{Gestern habe ich einen Bericht über die Gewalt in den Schulen gesehen.} — Parfait avec \all{haben} ; Partizip II de \all{sehen} : \all{gesehen} ; \all{in den Schulen} : datif pluriel.
\item \all{Die Arbeitslosigkeit hat schlechte Folgen für die Familien.} — « de mauvaises conséquences » : \all{schlechte Folgen} ; \all{für} + accusatif.
\end{enumerate}
\textbf{Points de vigilance :} ne pas traduire « parce que » par \all{denn} si la subordonnée suit le verbe : les deux sont possibles, mais avec \all{weil} le verbe va à la fin ; avec \all{denn} il reste en position 2 : \all{\ldots, denn sie haben keine Arbeit.}
\end{corrige}

\begin{corrige}[Exercice 6 — Transformation : denn et weil]
\begin{enumerate}
\item \all{Er lernt Deutsch, weil er in Deutschland studieren will.} — Le verbe conjugué \all{will} passe à la fin ; l'infinitif \all{studieren} le précède.
\item \all{Sie bleibt zu Hause, denn sie ist krank.} — Avec \all{denn}, retour à l'ordre direct : sujet + verbe en position 2.
\item \all{Wir helfen den Armen, weil sie kein Geld haben.} — \all{haben} rejeté à la fin.
\item \all{Die Kinder gehen nicht zur Schule, denn die Familie ist arm.} — Ordre direct après \all{denn}.
\end{enumerate}
\textbf{Points de vigilance :} \all{weil} = subordonnante (verbe à la fin) ; \all{denn} = coordinante (verbe en position 2). Le sens est identique ; seule la construction change.
\end{corrige}

\begin{corrige}[Exercice 7 — Texte à trous : le comparatif]
\all{In der Stadt ist das Leben \textbf{teurer} (teuer) als auf dem Land. Die Häuser sind \textbf{größer} (groß), aber viele Menschen sind \textbf{ärmer} (arm). Die Straßen sind nachts \textbf{gefährlicher} (gefährlich). Auf dem Land ist die Luft \textbf{besser} (gut) und die Menschen sind oft \textbf{freundlicher} (freundlich).}

\textbf{Justifications :}
\begin{itemize}
\item \all{teuer $\rightarrow$ teurer} : les adjectifs en \all{-er} perdent le \emph{e} au comparatif.
\item \all{groß $\rightarrow$ größer}, \all{arm $\rightarrow$ ärmer} : Umlaut sur \emph{a, o, u} pour les adjectifs monosyllabiques.
\item \all{gut $\rightarrow$ besser} : comparatif irrégulier, à apprendre par cœur.
\end{itemize}
\end{corrige}

\begin{corrige}[Exercice 8 — Appariement : mots et définitions]
\begin{center}
\begin{tabular}{|l|l|}\hline
\rowcolor{popblueL} Deutsch & Français \\ \hline
\all{die Arbeitslosigkeit} (1) & b) le chômage \\ \hline
\all{die Folge} (2) & a) la conséquence \\ \hline
\all{die Gewalt} (3) & d) la violence \\ \hline
\all{gefährlich} (4) & c) dangereux \\ \hline
\all{die Armut} (5) & e) la pauvreté \\ \hline
\end{tabular}
\end{center}
\textbf{Astuce mémo :} \all{die Arbeit} (le travail) + \all{-los} (sans) + \all{-igkeit} $\Rightarrow$ \all{die Arbeitslosigkeit} = « l'état d'être sans travail ».
\end{corrige}

\begin{corrige}[Exercice 9 — Type Bac : Compréhension écrite]
\textbf{Barème indicatif :} A : 5 $\times$ 2 pts = 10 pts ; B : 3 $\times$ 2 pts = 6 pts ; C : 2 $\times$ 2 pts = 4 pts. Total : 20 pts.

\textbf{A. Questions de compréhension}
\begin{enumerate}
\item \all{Es ist schwer, eine Arbeit zu finden, weil die Firmen oft Arbeiter mit Erfahrung wollen.} (Les entreprises veulent des travailleurs expérimentés.)
\item \all{Viele Jugendliche bleiben bei den Eltern, weil eine eigene Wohnung zu teuer ist.}
\item \all{Die Folgen sind Armut und Stress.} (ou : \all{Stress und Drogen.})
\item \all{In der Ausbildung lernen die Jugendlichen einen Beruf in einer Firma und in der Schule.}
\item \all{Die Experten sagen, dass Bildung der beste Weg gegen die Arbeitslosigkeit ist.}
\end{enumerate}

\textbf{B. Vrai ou faux ?}
\begin{enumerate}
\item \textbf{falsch} — \all{„Die Firmen wollen oft Arbeiter mit Erfahrung.“}
\item \textbf{falsch} — \all{„Eine eigene Wohnung ist zu teuer.“}
\item \textbf{richtig} — \all{„Viele finden danach eine gute Arbeit und verdienen mehr Geld als vorher.“}
\end{enumerate}

\textbf{C. Questions de langue}
\begin{enumerate}
\item Exemple de subordonnée avec \all{weil} : \all{„\ldots, weil junge Menschen ohne Arbeit kein Geld verdienen können.“} (ou : \all{„\ldots, weil eine eigene Wohnung zu teuer ist.“})
\item Deux comparatifs : \all{länger} ($\leftarrow$ \all{lang}) et \all{besser} ($\leftarrow$ \all{gut}) ; on accepte aussi \all{mehr} ($\leftarrow$ \all{viel}).
\end{enumerate}

\textbf{Points de vigilance :} répondre par des phrases complètes (sujet + verbe en position 2) ; recopier les justifications exactement, sans modifier le texte ; ne pas confondre \all{besser} (comparatif de \all{gut}) et \all{mehr} (comparatif de \all{viel}).
\end{corrige}

\begin{corrige}[Exercice 10 — Type Bac : Grammaire et conjugaison]
\textbf{Barème indicatif :} A : 6 pts ; B : 6 pts ; C : 8 pts. Total : 20 pts.

\textbf{A. La cause}
\begin{enumerate}
\item \all{Die Kinder hungern, \textbf{denn} die Familie \textbf{hat} kein Geld.} (ou \all{weil} : \all{\ldots, weil die Familie kein Geld hat.})
\item \all{Er nimmt keine Drogen, \textbf{denn} sie \textbf{sind} gefährlich.} (ou \all{weil sie gefährlich sind})
\item \all{Viele Jugendliche sind arbeitslos, \textbf{weil} es wenige Jobs \textbf{gibt}.} (ou \all{denn es gibt wenige Jobs})
\item \all{Wir bleiben zu Hause, \textbf{weil} es \textbf{regnet}.} (ou \all{denn es regnet})
\item \all{Sie lernt viel, \textbf{weil} sie eine gute Note \textbf{will}.} (ou \all{denn sie will eine gute Note})
\item \all{Die Stadt ist laut, \textbf{weil} es viele Autos \textbf{hat}.} (ou \all{denn es hat viele Autos})
\end{enumerate}
Les deux conjonctions sont acceptées partout, à condition que la place du verbe soit correcte : position 2 après \all{denn}, à la fin après \all{weil}.

\textbf{B. Le parfait}
\begin{enumerate}
\item \all{Der Junge \textbf{hat} keine Arbeit \textbf{gefunden}.} (\all{finden} : \all{haben} + \all{gefunden})
\item \all{Die Familie \textbf{ist} in die Stadt \textbf{gezogen}.} (déplacement $\Rightarrow$ \all{sein})
\item \all{Wir \textbf{haben} den armen Kindern \textbf{geholfen}.} (\all{helfen} + datif : \all{den Kindern})
\item \all{Sie \textbf{sind} lange im Krankenhaus \textbf{geblieben}.} (\all{bleiben} $\Rightarrow$ toujours \all{sein})
\item \all{Ich \textbf{habe} einen Film über Drogen \textbf{gesehen}.} (\all{sehen — sah — hat gesehen})
\item \all{Die Schüler \textbf{sind} nach Deutschland \textbf{gefahren}.} (\all{fahren} = déplacement $\Rightarrow$ \all{sein})
\end{enumerate}

\textbf{C. Version}
« Le chômage des jeunes est un problème social, parce que les jeunes sans travail ne peuvent pas gagner d'argent. Les conséquences sont dangereuses : la pauvreté, le stress et parfois aussi la drogue. Cette formation vaut mieux que rien, car les jeunes acquièrent de l'expérience et un diplôme. »

\textbf{Points de vigilance :} \all{verdienen können} = « pouvoir gagner » (double infinitif : \all{können} se traduit par « pouvoir ») ; \all{besser als nichts} = « mieux que rien » ; \all{denn} se traduit par « car ».
\end{corrige}

\begin{corrige}[Exercice 11 — Type Bac : Expression écrite guidée]
\textbf{Barème indicatif :} respect des 5 consignes : 5 pts ; correction grammaticale (place du verbe, parfait, comparatif) : 5 pts ; vocabulaire et orthographe (majuscules, Umlaute) : 5 pts ; cohérence et longueur : 5 pts. Total : 20 pts.

\textbf{Proposition de rédaction modèle :}

\all{Ein großes soziales Problem in Kamerun ist die Arbeitslosigkeit. Viele Jugendliche in Yaoundé und Douala finden keine Arbeit, weil es wenige Jobs gibt. Die Folgen sind gefährlich : Armut und manchmal auch Drogen. Das Leben in der Stadt ist teurer als auf dem Land. Letztes Jahr hat die Regierung ein Programm für die Jugendlichen gestartet. Ich denke : Mehr Schulen und mehr Ausbildung sind die beste Lösung gegen die Arbeitslosigkeit.}

(Traduction : « Un grand problème social au Cameroun est le chômage. Beaucoup de jeunes à Yaoundé et Douala ne trouvent pas de travail parce qu'il y a peu d'emplois. Les conséquences sont dangereuses : la pauvreté et parfois aussi la drogue. La vie en ville est plus chère qu'à la campagne. L'année dernière, le gouvernement a lancé un programme pour les jeunes. Je pense : plus d'écoles et plus de formations sont la meilleure solution contre le chômage. »)

\textbf{Vérification des consignes :}
\begin{itemize}
\item phrase avec \all{weil} : \all{\ldots, weil es wenige Jobs gibt.} (verbe \all{gibt} à la fin) ;
\item comparatif : \all{teurer als} ;
\item parfait : \all{hat \ldots\ gestartet} ;
\item lexique : \all{die Arbeitslosigkeit, die Folgen, gefährlich, die Armut, die Drogen} ;
\item solution finale : \all{Mehr Schulen und mehr Ausbildung sind die beste Lösung\ldots}
\end{itemize}

\textbf{Points de vigilance :} majuscule à tous les noms ; verbe en position 2 dans les phrases principales ; auxiliaire \all{haben} avec \all{starten} ; ne pas écrire \all{mehr gut} mais \all{besser} ; éviter les phrases trop longues : mieux vaut des phrases courtes et correctes.
\end{corrige}

\newpage

\coursec{Fiche bilan}

\begin{popfigure}
\begin{tikzpicture}[mindmap, grow cyclic, every node/.style={concept}, concept color=popblue!60, text=popdark, level 1/.append style={level distance=3.2cm, sibling angle=51}, level 2/.append style={level distance=1.9cm, sibling angle=45}]
\node[concept color=popblue!80, text=white, font=\bfseries] {Fléaux sociaux ALA03}
  child[concept color=poporange!70] { node {Lexique} 
      child { node[font=\scriptsize] {die Armut} }
      child { node[font=\scriptsize] {die Drogen} }
      child { node[font=\scriptsize] {die Gewalt} }
  }
  child[concept color=popgreen!70] { node {La cause}
      child { node[font=\scriptsize] {weil + verbe final} }
      child { node[font=\scriptsize] {denn + verbe 2e} }
  }
  child[concept color=poppurple!70] { node {Comparatif}
      child { node[font=\scriptsize] {größer als} }
      child { node[font=\scriptsize] {je ... desto} }
  }
  child[concept color=poppink!70] { node {Perfekt}
      child { node[font=\scriptsize] {haben/sein} }
      child { node[font=\scriptsize] {Partizip II} }
  }
  child[concept color=popgold!80] { node {Méthode}
      child { node[font=\scriptsize] {lire, repérer} }
      child { node[font=\scriptsize] {résumer, rédiger} }
  }
  child[concept color=popblue!50] { node {Texte}
      child { node[font=\scriptsize] {Soziale Probleme} }
      child { node[font=\scriptsize] {Folgen} }
  };
\end{tikzpicture}
\legende{Carte mentale de la leçon ALA03 : les fléaux sociaux et leurs conséquences}
\end{popfigure}

\begin{bilan}
\textbf{1. Lexique clé à connaître par cœur}
\begin{itemize}
  \item \all{die Armut} (--) : la pauvreté ; \all{die Arbeitslosigkeit} (--) : le chômage
  \item \all{die Drogen} (Pl.) : les drogues ; \all{die Gewalt} (--) : la violence
  \item \all{die Folge}, \all{-n} : la conséquence ; \all{die Ursache}, \all{-n} : la cause
  \item \all{gefährlich} : dangereux ; \all{arm} : pauvre ; \all{krank} : malade
  \item \all{das Problem}, \all{-e} : le problème ; \all{die Gesellschaft}, \all{-en} : la société
  \item \all{helfen} (hilft, half, hat geholfen) : aider ; \all{sterben} (stirbt, starb, ist gestorben) : mourir
\end{itemize}

\textbf{2. Grammaire : la cause}
\begin{center}
\begin{tabularx}{\linewidth}{|X|X|X|}\hline
\rowcolor{popblueL} Conjonction & Place du verbe & Exemple \\ \hline
\all{weil} (parce que) & verbe conjugué à la \textbf{fin} & \all{Er bleibt zu Hause, weil er krank ist.} \\ \hline
\all{denn} (car) & verbe en \textbf{2e position} & \all{Er bleibt zu Hause, denn er ist krank.} \\ \hline
\end{tabularx}
\end{center}

\textbf{3. Le comparatif}
\begin{itemize}
  \item Formation : adjectif + \all{-er} (souvent avec Umlaut : \all{arm} $\rightarrow$ \all{ärmer}, \all{jung} $\rightarrow$ \all{jünger}).
  \item Comparaison : \all{größer als} (plus grand que) ; égalité : \all{so ... wie} (aussi ... que).
  \item Irréguliers : \all{gut} $\rightarrow$ \all{besser} ; \all{viel} $\rightarrow$ \all{mehr} ; \all{gern} $\rightarrow$ \all{lieber}.
\end{itemize}

\textbf{4. Le parfait (Perfekt)}
\begin{itemize}
  \item Auxiliaire \all{haben} ou \all{sein} (mouvement, changement d'état) + participe II \textbf{à la fin}.
  \item Verbes faibles : \all{gemacht}, \all{gearbeitet} ; verbes forts : \all{gegangen}, \all{geschrieben}.
  \item Verbes à particule : \all{eingekauft}, \all{ferngesehen}.
\end{itemize}

\textbf{5. Méthode express (commentaire de texte)}
\begin{enumerate}
  \item Lire le texte deux fois et repérer le thème et l'idée générale.
  \item Souligner les mots-clés du champ lexical des fléaux sociaux.
  \item Répondre aux questions avec des phrases complètes (verbe en 2e position).
  \item Résumer en 3 à 5 phrases avec \all{weil} et le Perfekt.
\end{enumerate}
\end{bilan}

\begin{aretenir}[Checklist avant le Bac]
\begin{itemize}
  \item $\square$ Je connais l'article et le pluriel du vocabulaire : \all{die Armut}, \all{die Gewalt}, \all{die Folge}, \all{-n}.
  \item $\square$ Je sais que \all{weil} renvoie le verbe conjugué à la fin de la subordonnée.
  \item $\square$ Je sais que \all{denn} laisse le verbe en deuxième position.
  \item $\square$ Je forme le comparatif avec \all{-er} et je pense à l'Umlaut (\all{ärmer}, \all{größer}).
  \item $\square$ Je connais les comparatifs irréguliers : \all{besser}, \all{mehr}, \all{lieber}.
  \item $\square$ Je choisis le bon auxiliaire au Perfekt : \all{sein} pour le mouvement et le changement.
  \item $\square$ Je place le participe II à la fin de la phrase au Perfekt.
  \item $\square$ Je mets une majuscule à tous les noms allemands (\all{die Armut}, \all{das Problem}).
  \item $\square$ Je fais attention aux faux amis : \all{bekommen} = recevoir (et non « devenir »).
  \item $\square$ Je réponds aux questions de compréhension par des phrases complètes.
  \item $\square$ Je sais résumer un texte en quelques phrases avec des connecteurs de cause.
  \item $\square$ Je relis ma production : verbe en 2e position, accords, orthographe.
\end{itemize}
\end{aretenir}

\findecours

\end{document}
